\section{Full Proofs: The Lemma 3 Write-Up and the \texorpdfstring{$\nim{3}$}{*3} Families}\label{S-star3}
This section holds four notes. The first writes out Lemma~3 of Section~\ref{S-wall} (its Section 2b), which the other proofs cite, and proves a further $\nim{2}$ family with $n+2$ pawns. The second and third prove the whole-board $\nim{3}$ with $n+2$ pawns for even $n\ge8$ and odd $n\ge11$; the fourth proves the $\nim{3}$ with $n+1$ pawns for every $n\ge12$ (Section~\ref{sec:star2-rings}). Each is typeset without changes to its mathematics.
\subsection{Note 1: research/ring\_n\_plus\_2\_proof/PROOF.md}
\subsection*{$P_n=\star2$ for every even $n\ge 8$ (whole-board game): proof note}
\textbf{Summary}

\begin{enumerate}
\item Theorem: in the whole-board game on $C_n$ (race rule "count", e.p. on, an e.p.-capturable pawn is not passed, the game ends at the first passed pawn) $P_n=\star2$ for every even $n\ge8$; $P_n$ has $n+2$ pawns. "$=\star2$" has T50's meaning: $P_n+\star h$ and $\star2+\star h$ have the same outcome for every nim-heap (H8).
\item The proof is T50's (\texttt{research/ring\_family\_proof/proof.tex}): Lemmas 1-3 are quoted in section 2 with one repaired definition ("pattern file"), and the proof of Lemma 3 (ii)-(iii), which \texttt{proof.tex} gives in two terse lines, is written out in section 2b. The wall of $P_n$ is E-type ($W2/B4$) on files $6..n-1$; cut $c=12$, base case $n_0=18$.
\item Every hypothesis of Lemmas 2 and 3 holds for $P_n$ (section 3). T50's argument uses only the wall, the zone and the cut; the wall phase, its two ends and the zone location differ from $E_n$, and none of them enters (one end case, the double step of $W2$ on file $n-2$ at $n=18$, is L instead of P; H2).
\item Finite checks, by 3-5 programs each (section 4): C1' the reduced game $R_{18}$ has 813 nodes, no e.p. flag, zone $\{0,2,3,4,5,6\}$, and $G+\star2$ is a P-position; C2' $P_n+\star2$ is a P-position for $n=8,10,12,14,16$. \texttt{run\_all.sh} was rerun (32 s): its four logs equal the saved ones apart from timings. Two adversarial verifiers (section 8) rebuilt everything with two further programs and found no difference, for all even $n\le60$.
\item Status: PROVED [AI] for even $n\ge18$ given T50's Lemmas 1-3, the new text of 2b, and C1'; CHECKED-2 for $n=8..16$. The hand proof has been checked by AI models only; section 2b is new and has not yet been refereed by a second model. Scope: whole-board game only, race rule "count" (not "mover"), not cylcgt's component-sum model (192 of the 813 nodes of $R_{18}$ split into components).
\end{enumerate}
\subsubsection*{0. Claims and labels}
PROVED = by hand, given the lemmas (all hand proofs here are [AI]: checked by AI models only). CHECKED-2 = two or more independently written programs agree (the programs are listed in each row). COMPUTED-1 = one program. Programs: \textbf{A} = T50's \texttt{rg.c} (C, unchanged); \textbf{B} = \texttt{reduced\_cyl.py}, a rebuild on cylcgt's \texttt{pawns.moves/passed/\_child}; \textbf{C} = \texttt{check3/eng.py}, the T50 verifier's Python engine (written from the rules, independent of A and B); \textbf{D} = \texttt{verify\_logic/engine.py}, a fresh engine written from the rules by the logic verifier; \textbf{E} = \texttt{verify\_compute/myeng.py}, a fresh engine written from the \texttt{pawns.py} rules by the compute verifier (it shares only the position encoding with C, and was differentially tested against \texttt{pawns.moves/passed/promote\_distance}: about 61k positions, 0 mismatches). All five were written from the same rule description, not from each other's code; a shared misreading of the rules is excluded only by that differential test.

\par\noindent{\small\setlength{\tabcolsep}{3pt}\begin{tabular}{@{}>{\raggedright\arraybackslash}p{0.074\linewidth}>{\raggedright\arraybackslash}p{0.619\linewidth}>{\raggedright\arraybackslash}p{0.286\linewidth}@{}}
\toprule
id & claim & label \\
\midrule
L1-L3 & T50 Lemmas 1-3 (reduction, wall pawns, insertion), quoted in section 2 & PROVED in T50 (refereed by T51 and by the T50 verifier); not re-proved here, except Lemma 3 (ii)-(iii), written out in 2b \\
L3' & Lemma 3 (ii)-(iii) written out (2b) & PROVED, new text; corroborated by D and E (0 mismatches on the nodes, on node + 1..3 moves, and on exhaustive local tests); not yet refereed \\
H1 & the wall of $P_n$ is E-type; T50's definitions of wall and pattern file transfer (with the repaired definition of section 2) & PROVED (inspection) \\
H2 & Lemmas 2, 3 do not depend on the wall's phase or on its two ends & PROVED (re-reading of the proofs) \\
H3 & under Lemma 2's hypotheses $q$ has no capture and no e.p. capture & PROVED \\
H4 & the race rule enters only through the passed-pawn set of a child, which $\mathrm{ins}$ preserves & PROVED \\
H5 & $\mathrm{start}_n$ has no passed pawn and no flag; $\mathrm{ins}_{12}(\mathrm{start}_n)=\mathrm{start}_{n+2}$ & PROVED \\
H6 & cut and zone: files $8..17$ are pattern files in every node of $R_n$, $n\ge18$; the cut only enters through the hypotheses of Lemma 3 & PROVED given C1', by the induction of section 5 together with H7 \\
H7 & $R_{n+2}=\mathrm{ins}_{12}(R_n)$, with no flag and the same zone, $n\ge18$ & PROVED given L3, H5, C1' \\
H8 & for $h\ge3$ the sum $G+\star h$ is N (heap move to $G+\star2$); meaning of "$=\star2$"; the e.p. convention across heap moves is immaterial & PROVED \\
C1' & $R_{18}$: 813 nodes, 1643 K-edges, 1372 T-edges, no flag, zone $\{0,2,3,4,5,6\}$, $G+\star2$ is P & CHECKED-2 (A, B, C, D, E) \\
C2' & $P_n+\star2$ is P for $n=8,10,12,14,16$ & CHECKED-2 (A, C, D, E) \\
T18 & $P_n=\star2$ for even $n\ge18$ & PROVED [AI] given L1-L3, L3', C1' \\
T & $P_n=\star2$ for even $n\ge8$ & PROVED [AI] given C1', C2' (finite checks) \\
S1 & $R_n$ has 813 nodes and zone $\{0,2..6\}$ for every even $14\le n\le60$; the $\mathrm{ins}_{12}$ pump $18\to60$ (21 steps) & CHECKED-2 (A, C, E all $n\le60$; D up to $n=36$, pump up to $n=44$); redundant given H7 \\
S2 & unreduced $G+\star h$ has signature NNPN ($h\le3$; NNPNNNNN for $h\le7$) for every even $8\le n\le60$ & CHECKED-2 (A, C, E all $n\le60$; D for $n\le26$, 30, 36); not used beyond $n=16$ \\
S3 & cylcgt canonical value of $R_n$ is $\star2$ for $n=8..20$; whole-board cylcgt value of $P_8,P_{10}$ is $\star2$ & COMPUTED-1 (cylcgt arithmetic; a further route to C1', C2'; not used) \\
S4 & tests of Lemmas 2, 3 with the $P$-phase wall: random (C), exhaustive local with 2-4 outside files, node + 1..3 moves, random walks (D), Lemma 2 and 3 on all nodes for $n=18,20,22,26$ (E): 0 failures & CHECKED-2 (C, D, E); sanity only \\
\bottomrule
\end{tabular}}\par\smallskip
\par\noindent{\small\setlength{\tabcolsep}{3pt}\begin{tabular}{@{}>{\raggedright\arraybackslash}p{0.074\linewidth}>{\raggedright\arraybackslash}p{0.619\linewidth}>{\raggedright\arraybackslash}p{0.286\linewidth}@{}}
\toprule
id & claim & label \\
\midrule
S5 & scope: 192 of 813 nodes of $R_{18}$ split into components; 58 of 240 K-captures split & CHECKED-2 (C, E) \\
S6 & race-rule sensitivity: under "count" and "opponent" the signature is NNPNNNNN for $n=8..22$; under "mover" it is NRRRRRRR and $R_{18}$ has 553 nodes & COMPUTED-1 (E); informational \\
\bottomrule
\end{tabular}}\par\smallskip
\subsubsection*{1. Setting}
Model, T/L/P/K classes and reduced game: exactly as in \texttt{proof.tex} ("Model", lines 20-26). Terminal = at least one side has a passed pawn; non-terminal = neither has. Race rule "count" (\texttt{rg.c} \texttt{cls/pdist}, \texttt{pawns.\_child}; \texttt{RACE\_RULE = "count"} in \texttt{cylcgt/pawns.py}): if both have one after a move of $X$, the side with the smaller promotion distance wins, ties to $Y$ (who did not just move). The distance is the minimum over the side's passed pawns that are the front pawn of their file and not e.p.-capturable; $d=8-r-[r=2]$ (white, rank $r$), $d=r-1-[r=7]$ (black). $P_n$ ($n$ even $\ge8$), files $0..n-1$: $W5,\ B7,\ W4,\ W4B7,\ B6,\ W3B6$, then $(W2,B4)^{(n-6)/2}$. File $j\ge6$ holds $W2$ if $j$ is even, $B4$ if $j$ is odd. $R_n$ = reduced game of $G_n=P_n$ (nodes = positions with flag reachable by K-moves), $\mathrm{start}_n$ its root. $\mathrm{ins}_c$ as in Lemma 3.

\subsubsection*{2. T50's lemmas (statements quoted from proof.tex, Lemmas 1-3)}
The statements of Lemmas 1-3 below are quoted verbatim. The definition they rely on is not: \texttt{proof.tex} says "a pattern file is a wall file that still holds its starting pawn", which does not exclude extra pawns on the file, and Lemma 2's proof ("neighbours of $q$ are pawns of the other colour two ranks further on") needs the exact content. Throughout this note:

\textgreater{} \textbf{Pattern file.} A wall file whose content is \emph{exactly} its starting pawn and nothing else. The wall is files $3..n-1$ of $O_n$, $E_n$, resp. files $6..n-1$ of $P_n$: single pawns alternating $W3/B5$ (O) or $W2/B4$ (E), black two ranks above white. A \emph{wall pawn} is the pawn of a pattern file.

This is the reading T50's own tests use (they compare contents with equality: \texttt{lem.py}, \texttt{pump.py}).

\textbf{Lemma 1 (Reduction).} A side wins $G+\star h$ moving first iff it wins the reduced game moving first. \textbf{Lemma 2 (Wall pawns).} Let $c$ have no passed pawn, with an e.p. flag (if any) protecting a pawn of the player not to move, and let $q$ be a wall pawn of the player $Z$ to move, with pattern files $j-1,j,j+1$. (a) After any move of $q$, $Z$ has no passed pawn. (b) If $c$ has no flag, $Z=X$, and files $j-1,\dots,j+3$ (or $j-3,\dots,j+1$) are pattern files, every move of $q$ is L or P. \textbf{Lemma 3 (Insertion).} Let $p$ on $C_n$ have no passed pawn and no e.p. flag, and let files $c-4,\dots,c+5$ be pattern files. Let $\hat p=\mathrm{ins}_c(p)$ on $C_{n+2}$ be $p$ with two files, with the contents of files $c+1$ and $c$, inserted after file $c$ (the pattern continues). For each side, the children of its K-moves in $\hat p$ are the $\mathrm{ins}_c$-images of those in $p$, and it has a T-move in $\hat p$ iff in $p$.

\paragraph*{2b. Proof of Lemma 3, written out (new text)}
\texttt{proof.tex} proves Lemma 3 in six lines; its step (iii) cites (ii) for the replies of $Y$, although (ii) is about moves of $X$. What (iii) needs is the following, which is proved here from Lemma 2 and the view lemma V.

\emph{Notation.} $\mathrm{ins}_c(q)$ has files $q_0..q_c$, then $q_{c+1}$ and $q_c$ (the two \emph{new} files $c+1,c+2$), then the old files $q_{c+1},q_{c+2},\dots$ at places $c+3,c+4,\dots$ (old file $f\ge c+1$ moves to $f+2$). The pawns on old files are the \emph{image} pawns. The \emph{view} of a pawn is the contents of its own file and the two neighbouring files (the e.p. flag aside); a pawn's moves and its passed status depend on its view and the flag only. $\mathrm{ins}_c$ carries a flag along with its pawn.

\textbf{V (view lemma).} Let $q$ be any position on $C_n$. (V1) Every image pawn of $\mathrm{ins}_c(q)$ has the same view as its preimage in $q$ (the file left of old file $c+1$ is a copy of $q_c$, the file right of old file $c$ is a copy of $q_{c+1}$; nothing else changes). (V2) A new pawn has the view $(q_c,q_{c+1},q_c)$ or $(q_{c+1},q_c,q_{c+1})$. If files $c,c+1$ of $q$ are pattern files, these are views in which the pawn is blocked by both neighbours (pattern pawns of the other colour, two ranks further on), so new pawns are not passed. (V3) Hence, if files $c,c+1$ of $q$ are pattern files, the passed pawns of $\mathrm{ins}_c(q)$ are exactly the images of those of $q$, with the same colour, rank, own-file contents and e.p. protection. So $q$ and $\mathrm{ins}_c(q)$ are both non-terminal, or both terminal with the same winner (the race rule compares the same pawns, and the same side moved last). In particular $\hat p$ has no passed pawn.

\emph{Proof of Lemma 3.} Let $X$ be the side to move and $Y$ the other. Put $A=\{c-4,\dots,c+5\}$, a set of pattern files of $p$; then $\{c-4,\dots,c+7\}$ are pattern files of $\hat p$. (i) \emph{Central pawns.} A pawn of $X$ on a file $j\in\{c-3,\dots,c+4\}$ of $p$ has its three files in $A$ and a five-file window $j-1..j+3$ (for $j\le c+2$) or $j-3..j+1$ (for $j\ge c-1$) inside $A$. By Lemma 2(b) ($p$ has no passed pawn and no flag) every move of it is L or P, so it is neither T nor K. The same holds in $\hat p$ for the pawns on files $c-3,\dots,c+6$ (the new pawns included), whose windows lie inside the pattern arc $c-4..c+7$ of $\hat p$. (ii) \emph{Far moves of $X$.} Every other pawn of $X$ in $p$ stands on a file $\le c-4$ or $\ge c+5$; it corresponds to a pawn of $\hat p$ on a file $\le c-4$ or $\ge c+7$, with the same moves (V1). A move $\mu$ of it touches only files $\le c-3$ or $\ge c+4$ (a capture changes the neighbouring file; the file of an e.p. victim is the destination file; a double step puts the flag on its own file). So files $c-2,\dots,c+3$ of the child $c_\mu$ are unchanged, files $c,c+1$ are pattern files, and $\mathrm{ins}_c(c_\mu)$ is the child of the corresponding move $\hat\mu$. By V3, $c_\mu$ and $\mathrm{ins}_c(c_\mu)$ are both non-terminal or both terminal with the same winner, i.e. $\mu$ and $\hat\mu$ are both T, both L, or both non-terminal. (iii) \emph{Poison.} Let $c_\mu$ be non-terminal. $\mu$ is P iff $Y$ has a reply $\nu$ with $c_\mu\nu$ terminal and won by $Y$ (then $Y$ has a passed pawn). Files $c-2,\dots,c+3$ of $c_\mu$ hold the pattern, and files $c-2,\dots,c+5$ of $\mathrm{ins}_c(c_\mu)$ do. Call a pawn of $Y$ \emph{near} if it stands on a file in $c-1..c+2$ of $c_\mu$ (resp. in $c-1..c+4$ of $\mathrm{ins}_c(c_\mu)$, which contain the two new files), \emph{far} otherwise (files $\le c-2$ or $\ge c+3$, resp. $\le c-2$ or $\ge c+5$).

\begin{itemize}
\item \emph{Near replies.} A near pawn is a wall pawn whose files $j-1,j,j+1$ lie inside the pattern files just listed, and $c_\mu$ has no passed pawn and an e.p. flag (if any) protecting a pawn of $X$, the player not to move. By Lemma 2(a), after any near reply $Y$ has no passed pawn, so the result is not terminal-won by $Y$. This holds in $c_\mu$ and in $\mathrm{ins}_c(c_\mu)$ alike.
\item \emph{Far replies.} The far pawns of $c_\mu$ and of $\mathrm{ins}_c(c_\mu)$ correspond (V1), with the same moves. A far reply $\nu$ touches only files $\le c-1$ or $\ge c+2$ (a pawn on $c+3$ may capture onto $c+2$; one on $c-2$ onto $c-1$), so files $c,c+1$ of $c_\mu\nu$ are still pattern files and $\mathrm{ins}_c(c_\mu\nu)$ is the child of $\hat\nu$. By V3, $c_\mu\nu$ is terminal-won by $Y$ iff $\mathrm{ins}_c(c_\mu\nu)$ is. (The e.p. victim of an earlier far move $\mu$ sits on a file $\le c-4$ or $\ge c+5$, so e.p. replies are far replies.) Hence $Y$ has a winning terminal reply to $\mu$ iff to $\hat\mu$: $\mu$ is P iff $\hat\mu$ is. By (i)-(iii) the T-moves of $\hat p$ are the images of those of $p$ (so a T-move exists in $\hat p$ iff in $p$), and the K-moves of $\hat p$ are the images of the K-moves of $p$, with children $\mathrm{ins}_c(c_\mu)$. $\qquad\blacksquare$
\end{itemize}
\emph{Remark.} The argument never uses the phase of the wall (which colour stands on file $c$) or any property of the files outside $A$. It also shows that the only role of the cut is to guarantee that $A\subset$ pattern files (the hypothesis); it is not otherwise special (H6).

\subsubsection*{3. Hypotheses of Lemmas 2 and 3 for $P_n$}
\textbf{H1 (wall, pattern file).} With the definition of section 2, the wall of $P_n$ is files $6..n-1$, alternating $W2$ (even $j$) and $B4$ (odd $j$): type E, black pawn two ranks above the white one. In the nodes of $R_n$ ($n\ge14$) files $7..n-1$ always equal their start content, so the exact-content reading costs nothing: they are pattern files. \textbf{H2 (phase, ends).} $E_n$ has $W2$ on odd files and ends with $W2$ on file $n-1$; $P_n$ has $W2$ on even files, starts with $W2$ on file 6 (next to $W3B6$) and ends with $B4$ on file $n-1$ (next to $W5$ on file 0). Lemma 2(b) uses only: the capturer stands on $j+1$ (a $Y$ pawn), the $X$ pawn on $j+2$ is blocked by the $Y$ pawn on $j+3$ (rows $W2\to3$, $W2\to4$ e.p., $B4\to3$ of T50's table), or the mirror image with the window $j-3..j+1$; and, when $Y$ already has a passed pawn after the move, the move is L. Lemma 3 uses only V1-V3 and Lemma 2: the views of pawns are compared by content, and the new files carry the pattern for either phase (the inserted pair $(x_{c+1},x_c)$ keeps period 2). Neither uses a file's parity or a wall end. The content assumed is only on the arc $c-4..c+5$ (resp. $j-3..j+3$) of pattern files; everything outside the arc is arbitrary. The outside of the arc is one cyclic arc and is left unchanged by $\mathrm{ins}_{12}$. For $n=18$ the arc is $8..17$, it ends at file $n-1$, and the outside is files $0..7$; for $n>18$ the outside is files $18..n-1$ together with $0..7$ and contains both $n-1$ and $0$. Either way the wrap at the zone is harmless. \emph{An end case.} At $n=18$ the pawn $W2$ on file $16=n-2$ has the window $13..17$ on its left (the window on its right would need file 0, which holds $W5$). Its single step is P (the reply is the capture by $B4$ on file 15), but its double step is L, not P: it lands on rank 4, so $B4$ on file 17 has no white pawn ahead of it on files $16,17,0$ ($W5$ is behind it) and becomes passed. This is the branch "if $Y$ has a passed pawn, the move is L" of Lemma 2(b). In $p$ and in $\mathrm{ins}_{12}(p)$ this pawn lies on a central file ($16=c+4$, resp. $18=c+6$) and its moves are L or P in both, so nothing breaks. (Computed: this double step is L in all 813 nodes at $n=18$; for $n=20$ it is P.) \textbf{H3 ($q$ cannot capture).} Under Lemma 2's hypotheses files $j\pm1$ hold only the other colour's pattern pawn, two ranks further on, so the capture square is empty. An e.p. capture by $q=B4$ would need a white pawn that just stepped $2\to4$ on file $j\pm1$, but then that file is not a pattern file; $q=W2$ is on rank 2 and cannot capture e.p. (this fills the gap noted by the T51 referee: Lemma 2(b)'s table lists only forward steps). \textbf{H4 (race rule).} Lemma 2 never invokes the race (it shows $Z$ has no passed pawn). In Lemma 3, V3 gives that $c_\mu$ and $\mathrm{ins}_c(c_\mu)$ have the same passed pawns (same views; the new pawns stay blocked), with the same rank, colour, own-file occupancy and e.p. protection, hence the same class T/L/non-terminal under "count" (and under any rule that is a function of the passed pawns and of who moved). In (iii) the near replies never give $Y$ a passed pawn (Lemma 2(a)), so no race is won by $Y$. The proof itself is therefore not specific to $P_n$ or to "count"; the \emph{value} is (S6: under the rule "mover" the signature is NRRRRRRR and $P_n\ne\star2$). \textbf{H5 (start).} $\mathrm{start}_n$ has no passed pawn: $W5$(0) is blocked by $B7$(1); $W4$(2) by $B7$(1); $W4$(3) by $B7$(3); $W3$(5) by $B6$(4); $W2$(even $j$) by $B4$ on $j\pm1$ (file 6: $B4$(7)); $B7$(1) by $W5$(0); $B7$(3) by $W4$(2); $B6$(4) by $W4$(3); $B6$(5) by $W3$(5); $B4$(odd $j$) by $W2$ on $j-1$ (rank 2 \textless{} 4; for $j=n-1$ via $W2$(n-2), the $W5$ on file 0 is irrelevant). No flag. Hand check of the insertion: files $0..12$ are unchanged; the new files are copies of file 13 ($B4$) and file 12 ($W2$); the result alternates $W2$ (even) / $B4$ (odd) from file 6 on and has $n+2$ files, so $\mathrm{ins}_{12}(\mathrm{start}_n)=\mathrm{start}_{n+2}$ (also \texttt{start-{}\textgreater{}start True} in \texttt{logs/pump\_P\_fixc12.txt}). \textbf{H6 (zone, cut).} Zone $Z_0=\{0,2,3,4,5,6\}$ (file 1 = $B7$ never changes); the wall files $7..n-1$ are outside it. Admissible cuts need $c-4\ge7$ and $c+5\le n-1$, i.e. $11\le c\le n-6$; take $c=12$, files $8..17$. The arc does not wrap iff $n\ge18=n_0$ (T50's referee asked for this bound to be explicit). Unlike the zones of $O_n$, $E_n$, $Z_0$ lies left of $c$, does not contain file $n-1$, and is not shifted by $\mathrm{ins}_{12}$. $n_0=18$ is forced only by the radius 5 (smallest even $n$ with ten wall files $\ge7$): $R_n$ already has 813 nodes and the same zone for $n=14,16$, which are covered by C2'. \emph{The cut enters only through the hypotheses.} In every node of $R_n$ the files $7..n-1$ hold the period-2 wall, so inserting one period $(x_{c+1},x_c)$ after any file $c$ with $7\le c\le n-2$ gives the same position: $\mathrm{ins}_c(p)=\mathrm{ins}_{12}(p)$ for all nodes $p$ and all such $c$ (checked for all 813 nodes at $n=18$; the pump check $R_{n+2}=\mathrm{ins}_c(R_n)$ passes for every $c=7..n-2$ and fails only for $c=6$, which cuts the zone). The radius-5 arc and hence the range $11\le c\le n-6$ is only what Lemma 3 needs to be applicable; the pump at $c=12$ is not sensitive to the choice. No other property of $O_n$, $E_n$ is used in T50's proof: Lemma 1 is about the game, Lemmas 2, 3 about the wall, and the Theorem's proof about the base graph (C1), the cut and $\mathrm{ins}_c(\mathrm{start}_n)=\mathrm{start}_{n+2}$. The only places where the family enters are C1', C2' and H5/H6. \textbf{H8 (heaps).} \emph{Meaning.} As in T50 (model paragraph, Lemma 1, Theorem), "$P_n=\star2$" means: for every nim-heap $\star h$ ($h<256$ in T50's model, where a passed pawn is worth $\pm256$; the step below uses no bound) $P_n+\star h$ and $\star2+\star h$ have the same outcome, namely N for $h\ne2$ and P for $h=2$. It is the whole-board/heap-sum convention; it does not assert $P_n+X\equiv\star2+X$ for an arbitrary game $X$. \emph{Large heaps.} The programs compute the outcome of $G+\star h$ for $h=0..3$ (and $0..7$ in the logs): N, N, P, N (NNPNNNNN). For every $h\ge3$ no computation is needed: $G+\star h$ has the heap move $\star h\to\star2$ to $G+\star2$, a P-position, so $G+\star h$ is N (the start has no flag, so this heap move is an ordinary move of the sum). Thus the outcomes of $P_n+\star h$ and $\star2+\star h$ agree for all $h$. \emph{E.p. across a heap move.} \texttt{proof.tex} does not say whether a pending e.p. flag survives a heap move. It does not matter for the proof: the nodes of $R_n$ carry no flag, K-children never carry one, and the L/P arguments of Lemma 1 are direct replies that involve no heap move. For the unreduced searches (C2', S2) both readings ("keep", cylcgt's sum semantics; "expire", textbook chess) were run by E and give identical node counts and signatures for every even $n=8..60$.

\subsubsection*{4. Finite checks (rerun: \texttt{sh research/ring\_n\_plus\_2\_proof/run\_all.sh}, \texttt{run\_check3.sh}; logs in \texttt{logs/}; verifier reruns in \texttt{verify\_compute/}, \texttt{verify\_logic/})}
\textbf{C1' ($n_0=18$).} $R_{18}$ has 813 nodes, 1643 K-moves (all to distinct children) and 1372 T-moves (555 (node, side) pairs have one, so B's 2198 = 1643 + 555 distinct (node, side, target) triples); no node has an e.p. flag; no node has a passed pawn; every node agrees with $\mathrm{start}_{18}$ outside files $\{0,2,3,4,5,6\}$; solving the reduced game gives signature NNPN for $h=0..3$ (NNPNNNNN for $h\le7$), i.e. $G+\star2$ is a P-position. Sources: A \texttt{pump\_P\_8\_60.txt} (also nodes/zone/flags for every even $n\le60$: 813 and the same zone from $n=14$); B \texttt{reduced\_cyl\_P.txt} (node and edge sets identical to A for $n=8..22$; node strings include any \texttt{ep} token, so flags and zone agree); C \texttt{c3\_reduced\_P.txt} (813, 1643, 1372, flagged 0, touched files $\{0,2,3,4,5,6\}$, NNPN, $n=8..24$); D \texttt{verify\_logic/logs\_basic.txt}, \texttt{t\_counts.py} (the same numbers for every even $n=14..36$; 745, 1139, 889 nodes at $n=8,10,12$); E \texttt{verify\_compute/logs/b\_reduced\_8\_60.txt} (graph identical to C's, node by node, T-sets and K-sets, for all 27 even $n=8..60$, and equal to A's logs: \texttt{compare\_logs.txt}, 0 differences). Third route: \texttt{reduced\_value\_P.txt}, the canonical form of $R_n$ in cylcgt's CGT arithmetic is exactly $\star2$ for $n=8..20$. \textbf{C2' (n = 8..16).} Unreduced search for $G+\star h$ (no poison pruning): NNPN for $n=8,10,12,14,16$ in A (\texttt{unreduced\_P\_8\_60.txt}, 27 lines for $n=8..60$), in C (\texttt{c3\_unreduced\_P.txt}, $n=8..22$), in D (\texttt{verify\_logic/logs\_unreduced.txt}, $n=8..26,30,36$) and in E (\texttt{verify\_compute/logs/a\_outcomes\_*.txt}, all even $n=8..60$, both e.p. readings; for $n=8,10$ also an exhaustive solver that evaluates every child of every reachable state, 100851 and 687117 states). cylcgt whole-board value of $P_8$, $P_{10}$ is $\star2$ (\texttt{whole\_cyl\_P.txt}). \textbf{Insertion on the actual graphs (S1).} $\mathrm{ins}_{12}(R_n)=R_{n+2}$ (nodes onto nodes, K-edges, T-flags, T-counts) for the 21 steps $n=18..58$ in C and E (\texttt{c\_insertion\_8\_60.txt}, 42 of 42 ok), and for $n=18..44$ in D (\texttt{logs\_audit2\_c12\_c11.txt}, also $c=11$). \textbf{Lemma 3 and 2 on the actual nodes (sanity).} C: all nodes of $R_{18}$, $R_{20}$ and all admissible cuts: 4580 and 7832 node-side-cut checks, 0 mismatches (\texttt{c3\_lemma3\_nodes\_P.txt}); E re-did this for $n=18,20,22,26$ (4580, 7832, 11084, 17588 checks, 0 mismatches), and with the arc requirement dropped it finds 774 and 422 mismatches at $c=5,6$ ($n=18$), so the test can fail. D: Lemma 2(a) and (b) on every move of every wall pawn with pattern neighbours (11382, 13821, 21138 moves at $n=18,20,26$; 0 captures, 0 passed-pawn violations, no T or K move); Lemma 3 on every node plus 1, 2 or 3 arbitrary moves (up to 32807 cut checks at $n=18$, 78620 at $n=20$, 76314 at $n=24$), on exhaustive local positions with 2-4 arbitrary outside files (1243, 49149, 109748 positions per phase) and on random walks: 0 mismatches; arcs of 2 and 4 files fail (897 and 99 mismatches), arcs of 6 files and longer pass, so radius 5 is conservative. C random tests with the $P$ phase: Lemma 2 (a) 9220 moves, (b) 6992 moves, 0 failures; Lemma 3 1857 accepted positions, 0 mismatches; with the arc shortened by 3 files the same test finds a mismatch (1 of 42 accepted).

\subsubsection*{5. Proof of the Theorem}
Let $n\ge18$ be even and let $I_n$ say: no node of $R_n$ has a flag, and every node agrees with $\mathrm{start}_n$ outside $Z_0$. $I_{18}$ is C1'. \emph{Step.} Assume $I_n$. Every node $p$ of $R_n$ is non-terminal (no passed pawn), has no flag, and its files $8..17$ are wall files $\ge7$ outside $Z_0$, so they hold their start content: they are pattern files (H1, H6). Lemma 3 with $c=12$ applies to every node: for each side the K-children of $\mathrm{ins}_{12}(p)$ are the $\mathrm{ins}_{12}$-images of those of $p$, and the T-flags agree. As $\mathrm{ins}_{12}(\mathrm{start}_n)=\mathrm{start}_{n+2}$ (H5), induction on the distance from the root shows that the nodes, K-edges and T-flags of $R_{n+2}$ are the $\mathrm{ins}_{12}$-images of those of $R_n$. Images have no flag, and since $Z_0\subset[0,6]$ lies left of the cut, $\mathrm{ins}_{12}(p)$ agrees with $\mathrm{start}_{n+2}$ outside $Z_0$: $I_{n+2}$ holds. So $R_n\cong R_{18}$ for all even $n\ge18$ (nodes, edges, T-flags). By C1', $G+\star2$ is a P-position of the reduced game; by Lemma 1 it is a P-position of $P_n+\star2$. The signature of $G+\star h$ is NNP for $h=0,1,2$ and, for every $h\ge3$, $G+\star h$ is N by the heap move to $G+\star2$ (H8); this is the outcome of $\star2+\star h$ for every $h$, i.e. $P_n=\star2$ (H8 for the meaning). For $n=8,\dots,16$ use C2' (the same H8 step applies). $\qquad\blacksquare$ (T18: PROVED [AI] given L1-L3, L3', C1'; T: with C2'.)

\subsubsection*{6. Gaps and caveats}
\begin{enumerate}
\item \textbf{No gap in the hypotheses.} Nothing in T50's Lemmas 1-3 or in the Theorem's induction uses a property of $O_n$/$E_n$ beyond the wall; the differences (phase, ends, zone location, $c=12$, $n_0=18$) were checked one by one (H2, H5, H6). The only repairs are precisions: the definition of "pattern file" (section 2), the explicit bound $n\ge n_0$ (H6), the no-capture remark (H3), the large-heap step and the meaning of "$=\star2$" (H8), and the written-out step (iii) (2b).
\item \textbf{Lemmas 1-3 are cited, not re-proved, except Lemma 3 (ii)-(iii).} Their hand proofs are T50's, refereed as sound in \texttt{ring\_family\_proof\_check/REPORT.md} and by the T50 verifier (ACCEPT\_WITH\_FIXES: presentation only). The race rule is not defined in \texttt{proof.tex}; H4 states it. Lemma 3 (iii) was terse in \texttt{proof.tex}; 2b replaces it. \textbf{2b is new text, so the hand proof of T18 rests on \texttt{proof.tex}'s Lemmas 1-2, on 2b, and on C1'; it has been checked by AI models only and 2b has not yet been refereed by a model other than the author.} The correspondence that 2b proves was tested by programs D and E with 0 mismatches (section 4), which is evidence but not proof. A non-Claude referee (the T53 brief) should receive 2b.
\item \textbf{C1', C2' are computer facts, not hand proofs.} The theorem for $n\ge18$ is "proved modulo C1'". Programs A, B, C, D, E agree on C1' and A, C, D, E on C2' (B is cylcgt's own code). All were written from one rule description (A, C, D, E share the model but not code; E shares only the position encoding with C); a misreading of the rules common to all of them would not be caught, apart from E's differential test against \texttt{pawns.py} (0 mismatches on about 61k positions, 690k child verdicts, 2.9k e.p. captures). The unreduced searches short-circuit on the first winning move; the version that evaluates every child was run only for $n=8,10$ (E); for $n\ge12$ the results rest on short-circuit solvers (C and E agree node for node).
\item \textbf{Scope.} Whole-board game only. In $R_{18}$, 192 of 813 nodes split into $\ge2$ components and 58 of 240 K-captures split the ring, so Lemma 1 is not available in cylcgt's default component-sum model; $P_n=\star2$ there is not covered by this note. \textbf{Race rule.} The result is for the rule "count" (the project's rule, \texttt{RACE\_RULE = "count"}); "opponent" gives the same signature on a different graph (881 nodes), but under "mover" (the mover wins every race) $P_n$ is not $\star2$ (NRRRRRRR, 553 nodes). The proof's lemmas are rule-independent (H4); C1' is not.
\item \textbf{Redundant evidence.} S1, S2 are now CHECKED-2 up to $n=60$ (A, C, E) and are not used beyond what H7 gives. S3 (cylcgt values) covers $n\le20$ only and is COMPUTED-1: \texttt{whole\_cyl\_P.txt} is copied from the orchestrating session's run (\texttt{scratchpad/np2/whole\_P.txt}) and its run for $n=12$ had not finished, so the whole-board cylcgt route is not a second program for C2' beyond $n=10$ (C2' stands on A, C, D, E).
\item $n=8,10,12$ have different graphs (745, 1139, 889 nodes; zones contain files 7 or 9), so they are not instances of the pump; only $n\ge14$ share $R_{18}$.
\end{enumerate}
\subsubsection*{7. Reproduce (from the project root)}
\texttt{sh research/ring\_n\_plus\_2\_proof/run\_all.sh} (32 s, A and B, writes the four original logs); \texttt{sh research/ring\_n\_plus\_2\_proof/run\_check3.sh} (about 2 min, C and cylcgt values, writes \texttt{c3\_*.txt}, \texttt{reduced\_value\_P.txt}). Verifier reruns: \texttt{sh research/ring\_n\_plus\_2\_proof/verify\_compute/run\_all.sh} (about 20 min, writes \texttt{verify\_compute/logs/}), and in \texttt{verify\_logic/} \texttt{python t\_basic.py 14 16 18 20 22}, \texttt{t\_unreduced.py}, \texttt{audit1\_hyp.py}, \texttt{audit2\_pump.py} etc. (logs next to the scripts). One process at a time, under \texttt{nice -{}n 15}; no git; nothing outside this directory edited.

\subsubsection*{8. Verification}
Two adversarial verifiers judged this note in separate lenses, then every problem they raised was re-checked against this file, \texttt{proof.tex}, their scripts and logs, and this section and the edits above were made. Their scripts and logs are in \texttt{verify\_logic/} and \texttt{verify\_compute/}.

\textbf{Verdicts.} Logic lens: ACCEPT\_WITH\_FIXES (could not refute that T50's Lemmas 1-3 apply to $P_n$; the fixes are wording and one missing line). Compute lens: ACCEPT (no difference found).

\textbf{What the logic verifier recomputed} (\texttt{verify\_logic/}, a fresh engine D written from the rules). C1': 813 nodes, 1643 K-moves, 1372 T-moves, 555 (node, side) pairs with a T-move, no flagged node, zone $\{0,2,3,4,5,6\}$, signature NNPN, for every even $n=14..36$; 745, 1139, 889 nodes at $n=8,10,12$. Hypotheses of Lemmas 2, 3 on all 813 nodes at $n=18,20,24$ (0 violations; K-moves come only from files $\{0,2,3,4,5,6\}$, T-moves only from $\{0,2,3,4,5\}$). Lemma 2(a) and (b) on 11382, 13821, 21138 moves (0 failures). $R_{n+2}=\mathrm{ins}_{12}(R_n)$ for $n=18..44$ (also for $c=11$). Lemma 3 on every node plus 1-3 moves, on exhaustive local positions with a menu of outside files in both phases, and on random walks: 0 mismatches; a sensitivity control (arcs of 2 and 4 files) fails as it should. E.p. flags: 5450 moves from the nodes at $n=18$ create a flag, 3562 are P and 1888 are L, none is K or T. Unreduced search without poison pruning: NNPN for $n=8..26,30,36$.

\textbf{What the compute verifier recomputed} (\texttt{verify\_compute/}, E written from the \texttt{pawns.py} rules, plus unchanged copies of C). (a) The full game $G+\star h$, $h=0..7$: NNPNNNNN for every even $n=8..60$ in three solvers (C's, E "keep", E "expire"), identical node counts; an exhaustive no-early-exit solver for $n=8,10$. (b) The reduced graph for every even $n=8..60$ in C and E, identical (813 nodes from $n=14$, K 1643, T 1372, flagged 0, no passed pawn in any node, zone $\{0,2,3,4,5,6\}$). (c) $\mathrm{ins}_{12}(R_n)=R_{n+2}$ for $n=18..58$ in both engines, and for every cut $c=7..n-2$ (fails only for $c=6$). Lemma 2 on the nodes at $n=18,22,26$ (0 violations). Lemma 3 on the nodes at $n=18,20,22,26$ (0 mismatches). A differential test of E, C and \texttt{pawns.py} (about 61k positions, 0 mismatches). Calibration of E on T50's families (O, E, F). The side counts 555, 427, 192, 240, 58 of this note. A comparison with the saved logs of A and B: 27 of 27 values of $n$ identical for the reduced graphs and the unreduced signatures, 21 of 21 pump steps, 0 differences. Race rules: NNPNNNNN under "count" and "opponent", NRRRRRRR under "mover".

\textbf{My own re-check} (a spot run of D's engine under \texttt{nice -{}n 15}, one process, not the 20-minute full rerun): $R_{18}$ has 813 nodes and signature NNPNNNNNNN for $h=0..9$; $\mathrm{ins}_{12}(R_{18})=R_{20}$ (813 nodes) and $\mathrm{ins}_{12}(\mathrm{start}_{18})=\mathrm{start}_{20}$; $\mathrm{ins}_c(p)$ is the same position for every $c=7..n-2$ and all 813 nodes; the moves of $W2$ on file 16 at $n=18$ are P (single step, 813 nodes) and L (double step, 813 nodes). Everything else was checked against the verifiers' logs and scripts, the numbers in this note against \texttt{compare\_logs.txt}, \texttt{c\_insertion\_8\_60.txt}, \texttt{b\_reduced\_8\_60.txt}, \texttt{a\_outcomes\_*.txt}, \texttt{race\_rules.txt}, \texttt{logs\_basic.txt}, \texttt{logs\_audit*.txt}.

\textbf{Problems raised, and what was done.}

\par\noindent{\small\setlength{\tabcolsep}{3pt}\begin{tabular}{@{}>{\raggedright\arraybackslash}p{0.072\linewidth}>{\raggedright\arraybackslash}p{0.412\linewidth}>{\raggedright\arraybackslash}p{0.072\linewidth}>{\raggedright\arraybackslash}p{0.434\linewidth}@{}}
\toprule
\# & problem (source) & confirmed? & fix \\
\midrule
1 & H2 said the outside of the arc contains files $n-1$ and $0$; false at $n=18$, where the arc $8..17$ contains file 17 (logic) & yes (harmless) & H2 rewritten: at $n=18$ the outside is files $0..7$, for $n>18$ it contains $n-1$ and $0$; one cyclic arc in both cases \\
2 & $P_n=\star2$ concluded from $h\le3$ (logs $h\le7$) only; no step for larger $h$; "G+\textbackslash{}\textasteriskcentered{}2 is P implies G=\textbackslash{}\textasteriskcentered{}2" is a convention (logic) & yes (missing line) & new H8 and a sentence in section 5: $G+\star h$ is N for $h\ge3$ by the heap move to $G+\star2$; meaning of "$=\star2$" stated; the e.p. convention across heap moves shown immaterial \\
3 & section 2 said "verbatim" but T50's "pattern file" definition is insufficient; the repair was only in H1 (logic) & yes & section 2 now quotes the statements verbatim and states the repaired definition (exact content) before them; "verbatim" no longer covers the definition \\
4 & Lemma 3 (iii) is terse and cites (ii) for $Y$'s replies, which is about $X$'s moves; "PROVED" rests on T50's terse proof (logic) & yes & section 2b: the proof of Lemma 3 written out with the view lemma V1-V3, a near/far split of $Y$'s replies, and the case of replies touching $c-1$, $c+2$; the label and summary now say that 2b is new and not yet refereed by a second model \\
5 & H6 does not say that the cut enters only through the hypotheses; $\mathrm{ins}_c(p)$ is the same for all $c\ge7$ (logic) & yes (clarification) & H6: $\mathrm{ins}_c(p)=\mathrm{ins}_{12}(p)$ for all nodes and $c=7..n-2$ (checked); only the radius-5 hypothesis restricts $c$ \\
6 & at $n=18$ the double step of $W2$ on file 16 is L, not P, since $B4$ on file 17 becomes passed; the note said the far end "does not enter" (logic) & yes (re-computed: L in all 813 nodes) & H2 "An end case": the L branch of Lemma 2(b) covers it; the move is L or P in $p$ and in $\mathrm{ins}_{12}(p)$ \\
7 & status labels: S1/S2 above $n=28$ / $n=22$ rest on A alone; S3 copied/unfinished; the summary must not suggest two programs everywhere (logic) and S1/S2 can be upgraded (compute) & yes & table: S1, S2 CHECKED-2 up to $n=60$ (A, C, E; D partly), C1' by 5 programs, C2' by 4, S3 COMPUTED-1 with the reason; summary item 4 says "3-5 programs each" \\
8 & the result is specific to the race rule "count" (and "opponent"); under "mover" $P_n\ne\star2$ (compute) & yes (informational) & S6, H4 and section 6.4 state it; the claim below names the rule \\
9 & independence caveat: A, C, D, E were written from the same rules (compute) & yes & section 0 and 6.3 say so and cite E's differential test against \texttt{pawns.py} \\
10 & no-early-exit solver was run only for $n=8,10$ (compute) & yes (caveat) & section 4 and 6.3 say so \\
11 & the e.p. flag across a heap move is not specified in \texttt{proof.tex} (compute) & yes (moot) & H8: immaterial for the proof; both readings give identical results for $n=8..60$ \\
\bottomrule
\end{tabular}}\par\smallskip
No problem was rejected. None changes the logic of the proof: problems 1, 5, 6 are statements about details that were already correct in effect, problem 2 adds a one-line step that follows from the heap move, problems 3 and 4 repair a definition and a terse argument, and problems 7-11 are labels and caveats. Section 2b is the only new mathematics, and the only part of the note that has not been through a referee.

\textbf{Final verdict: ACCEPT\_WITH\_FIXES (fixes applied).} Reasons: (1) the logic verifier could not refute the transfer of T50's lemmas, and found the hand part sound after the repairs of problems 2-4; (2) the compute verifier found no difference between two fresh engines, T50's engine, \texttt{rg.c} and \texttt{cylcgt} for any even $n\le60$, for the reduced graph, the pump, and the unreduced signatures; (3) the remaining weaknesses are named, not hidden: C1' and C2' are finite computations, 2b has not been refereed by a second model, and the result depends on the race rule.

\textbf{Sentence the paper may use.} In the whole-board game on the cylinder (race rule "count", en passant on), the wall ring $P_n$ of $n+2$ pawns (files $W5, B7, W4, W4B7, B6, W3B6$ followed by $(W2,B4)$ pairs) is $\star2$ for every even $n\ge8$, meaning that $P_n+\star h$ and $\star2+\star h$ have the same outcome for every nim-heap $\star h$ (of size below 256, the range of the model): for $n\ge18$ this follows from T50's insertion argument adapted by hand [AI] and one finite check (the 813-node reduced game at $n=18$), and for $n=8,\dots,16$ from direct finite searches, all confirmed by independently written programs [C2].

Do not add: that this holds in the component-sum model of cylcgt, for another race rule, or for odd $n$ (this note is about even $n$ only).


\subsection{Note 2: research/ring\_n\_plus\_2\_proof/PROOF\_Q.md (even $n$)}
\subsection*{The whole-board Q family is $\star3$ for every even $n\ge8$}
\subsubsection*{Summary}
\begin{enumerate}
\item \textbf{PROVED [AI], conditional on finite checks:} the insertion proof gives $Q_n=\star3$ for even $n\ge22$.
\item \textbf{CHECKED-2:} the $Q_{22}$ reduced graph has 7,968 nodes and signature NNNPN; two engines agree.
\item \textbf{CHECKED-2:} full-game signatures for even $n=8,\ldots,24$ are NNNP for $h=0,1,2,3$.
\item \textbf{COMPUTED-1:} an independent full-game run also checks $h=4$ and Lemma 2(b) on every $Q_{22}$ node.
\item Scope: whole-board game, race rule \texttt{count}, en passant on, and the e.p.-aware passed-pawn test.
\end{enumerate}
\subsubsection*{0. Claims and labels}
\textbf{PROVED} means a hand argument conditional on cited lemmas and finite checks; \textbf{CHECKED-2} means two independently written programs agree; \textbf{COMPUTED-1} means one program. The programs share the written rules, not implementation code.

\par\noindent{\small\setlength{\tabcolsep}{3pt}\begin{tabular}{@{}>{\raggedright\arraybackslash}p{0.382\linewidth}>{\raggedright\arraybackslash}p{0.558\linewidth}@{}}
\toprule
Claim & Status and evidence \\
\midrule
T50 Lemmas 1–3 and the insertion detail in \texttt{PROOF.md} §2b & \textbf{PROVED} in T50; the generic insertion argument is reused here. \\
H1–H8 and the induction below & \textbf{PROVED [AI]} conditional on C1 and on the T50 lemmas. \\
C1: $R_{22}$ graph, zone, flags, and reduced signature & \textbf{CHECKED-2}: A (\texttt{rg.c}, saved \texttt{pump\_Q\_8\_60.txt}) and E (\texttt{verify\_compute/myeng.py}, rerun in \texttt{q/check\_wall\_q22.py}); E also checked no passer at every node. \\
C2: unreduced signature NNNP, $h=0..3$, even $n=8..24$ & \textbf{CHECKED-2}: A (\texttt{unreduced\_Q\_8\_60.txt}) and F (\texttt{q/independent.py}). \\
F's $h=4$ result in both heap/e.p. interpretations; P8 control & \textbf{COMPUTED-1}. \\
Lemma 2(a),(b) on actual $R_{22}$ nodes & \textbf{COMPUTED-1}: E, 111,552 wall-pawn moves, zero failures. \\
\bottomrule
\end{tabular}}\par\smallskip
A is the saved C whole-board engine; E is the independently written Python engine in \texttt{verify\_compute/myeng.py}; F has its own move generation and whole-board recursion and calls only \texttt{cylcgt.pawns.passed} for the passed-pawn predicate. B, \texttt{reduced\_cyl.py}, rebuilds reduced graphs with \texttt{cylcgt}; its saved Q comparison stops at $n=8,10,12$, so it is \textbf{not} evidence for C1 at $n=22$ or for unreduced signatures. The pump log beyond $n=22$ is a one-program check and is not used in the proof.

\subsubsection*{1. Model and lemmas}
The game is the whole-board normal-play game on $C_n$: White is Left, Black is Right, and a side with no move loses. Race rule is \texttt{count}; en passant is on, and an e.p.-capturable pawn is not passed. A passer ends the game; if both sides have one, compare their shortest promotion distances, with a tie won by the side that did not just move. The terminal score convention is $\pm256$, so heaps considered in the value statement have $h<256$.

For even $n\ge8$, the files of $Q_n$ are \textbackslash{}[ W4,\textbackslash{} B7,\textbackslash{} (B4,W2)\textasciicircum{}\{(n-8)/2\},\textbackslash{} W3B6,\textbackslash{} W3,\textbackslash{} W2B5,\textbackslash{} B5,\textbackslash{} W2,\textbackslash{} B4. \textbackslash{}] Thus $Q_8$ is exactly \texttt{W4 B7 W3B6 W3 W2B5 B5 W2 B4} (the repeated block is empty), and every $Q_n$ has $n+2$ pawns. A \textbf{pattern file} is a file whose content is exactly its starting pawn and nothing else.

Use T50's reduction and wall-pawn lemmas as stated in \texttt{research/ring\_family\_proof/proof.tex}; the exact-content pattern-file reading and the written-out insertion reply argument are in \texttt{PROOF.md} §§2–2b. Lemma 2(b) says that, with no flag and no passer at the position, a moving wall pawn with a five-file pattern window has only L or P moves. The Q wall has the opposite phase from P: W2 is on odd files and B4 on even files. The same local proof applies: W2 advances $2\to3$ or $2\to4$ (the latter is e.p.-capturable), and B4 advances $4\to3$; the adjacent opposite-colour wall pawn captures the advance, creating a passer unless the move is already L. The X pawn on $j+2$ remains blocked by Y's pawn on $j+3$; the reversed window is its mirror. Lemma 2(a) excludes a passer for the mover after every such advance.

For Lemma 3 at $c=9$, the required pattern arc is $5..14$. In $p$, files $6..13$ have five-file windows inside that arc; after insertion, files $6..15$ do inside the extended arc $5..16$. As in \texttt{PROOF.md} §2b, central moves are L/P, far moves preserve the same three-file views and passed-pawn set, far replies correspond, and near wall replies cannot give Y a passer by Lemma 2(a). Hence T-flags and K-children correspond under insertion. No second or “symmetric” arc is used.

\subsubsection*{2. Q hypotheses H1–H8}
\begin{enumerate}
\item \textbf{H1 — Wall and pattern.} The periodic wall is files $2..n-7$: B4 on even files and W2 on odd files, with black two ranks above white. Pattern files require exact single-pawn content. The base graph has files $4..14$ fixed in every node; insertion preserves this, in particular $5..14$.
\item \textbf{H2 — Phase and the two relevant ends.} The ten-file arc for $c=9$ is $5..14$. At the two ends of the invariant pattern run, the applicable Lemma 2(b) windows are $4..8$ at $q$ on file 5 and $10..14$ at $q$ on file 13. On all actual $R_{22}$ nodes these windows hold; both cases and every intervening wall move were checked explicitly (C1). Lemma 3 invokes Lemma 2(b) only on central files $6..13$ (and $6..15$ after insertion), not at the physical wall endpoints 2 or $n-7$, where its five-file premise need not hold.
\item \textbf{H3 — No wall captures.} With pattern neighbors, a W2/B4 pawn's diagonal target on rank 3 is empty; B4 cannot e.p.-capture a W2 still on rank 2. Thus the wall moves in Lemma 2 are advances only. The lemma's L case covers any move after which Y already has a passer.
\item \textbf{H4 — Race and passed test.} This theorem uses \texttt{count} and the e.p.-aware passed test. The wall lemma gives no passer to the mover; insertion preserves passer identities, ranks, and e.p. protection on far moves, so it preserves the race result.
\item \textbf{H5 — Root and insertion.} The root has no passer: W4(0) is blocked by B7(1); B7(1) by W4(0); wall W2/B4 pawns by their adjacent opposite-colour wall pawns; W3B6, W3, W2B5, B5, and the closing W2/B4 pair by same-file or adjacent suffix pawns. It has no flag. For $n\ge22$, file 9 is W2 and file 10 is B4; inserting B4,W2 after file 9 extends the alternating block, so $\mathrm{ins}_9(\mathrm{start}_n)=\mathrm{start}_{n+2}$.
\item \textbf{H6 — Base zone and cut.} Take $n_0=22$, $c=9$. The base zone is $Z_{22}=\{0,1,2,3\}\cup\{15,16,17,18,19\}\cup\{21\}$. Files $5..14$ lie outside it and are exact pattern files. The bound is the first even $n$ for which $c+5=14\le n-7$.
\item \textbf{H7 — Inductive invariant.} Let $Z_n=\{0,1,2,3\}\cup\{n-7,\ldots,n-3\}\cup\{n-1\}$. Every node of $R_n$ has no passer or flag and agrees with its root outside $Z_n$. C1 gives this at $n=22$. If it holds at $n$, Lemmas 2–3 apply on $5..14$; H5 maps the root, and insertion maps every K-edge and T-flag. Since $\mathrm{ins}_9(Z_n)=Z_{n+2}$, the invariant holds at $n+2$ and $R_{n+2}$ is isomorphic to $R_n$ with T-flags.
\item \textbf{H8 — Meaning of $Q_n=\star3$.} It means that for every nim heap $\star h$, $h<256$, $Q_n+\star h$ and $\star3+\star h$ have the same outcome; it does not claim equivalence against arbitrary addends. It suffices to show NNNP at $h=0,1,2,3$: for every $h>3$, the move $\star h\to\star3$ reaches the P-position $Q_n+\star3$, so the sum is N, as is $\star3+\star h$. The root has no flag; the reduced nodes and K-children have none, so the heap-move e.p. convention does not affect this argument.
\end{enumerate}
\subsubsection*{3. Finite checks and theorem}
\textbf{C1 (CHECKED-2).} $R_{22}$ has 7,968 nodes, 21,063 K-edges, 17,784 T-edges, no flags, and zone $Z_{22}$; A and E agree on the graph and its NNNPN signature for $h=0..4$. E checked that no node has a passer and checked all 7,968 nodes and 71,712 q/node positions: 111,552 moves, zero Lemma 2(a) violations and zero Lemma 2(b) T/K violations. In particular, file-5 moves classify L/P as 5,691/10,245 and file-13 moves as 414/15,522. These end counts make the two boundary-window checks explicit.

\textbf{C2 (CHECKED-2 / COMPUTED-1).} A's saved unreduced log and F agree on NNNP for every even $n=8..24$. F additionally computes $h=4$ and both e.p.-across-heap readings (\texttt{keep}, \texttt{expire}): NNNPN for each of those nine sizes. The control $P_8$ is NNPNN in both readings, so it is not $\star3$ (its P is at $h=2$). The $h=4$ row and the control are COMPUTED-1 by F.

\textbf{Theorem (PROVED [AI], conditional on C1–C2 and T50 Lemmas 1–3).} H7 gives the same reduced game and T-flags as $R_{22}$ for every even $n\ge22$, so C1 and Lemma 1 give NNNP. H8 completes the heap comparison. For $8\le n<22$, C2 gives NNNP directly and H8 again completes it. Therefore $Q_n=\star3$ for every even $n\ge8$ in the stated whole-board model.

\textbf{Reproduction, from the project root:}

\begin{small}\begin{verbatim}
PYTHONPATH=. .venv/bin/python -u research/ring_n_plus_2_check/q/check_wall_q22.py
/usr/bin/time -l nice -n 15 env PYTHONPATH=. .venv/bin/python -u research/ring_n_plus_2_check/q/independent.py
\end{verbatim}\end{small}
The cited A outputs are the pre-existing \texttt{research/ring\_n\_plus\_2\_proof/logs/pump\_Q\_8\_60.txt} and \texttt{unreduced\_Q\_8\_60.txt}; their original reproduction command is \texttt{sh research/ring\_n\_plus\_2\_proof/run\_all.sh} (not run here because it writes outside \texttt{q/}). \texttt{reduced\_cyl\_Q.txt} is only a reduced-graph comparison for n=8,10,12 and is not used for the claims above.

\subsubsection*{4. Referee (T62)}
\textbf{Verdict: SOUND}, conditional on the finite computations C1–C2 and the cited T50 lemmas. No transfer gap remains for the even Q family at cut $c=9$ and base $n_0=22$.

\textbf{Recomputed:} the Q22 reduced graph and signature with E; Lemma 2(a),(b) on all 7,968 actual nodes, including both boundary windows; the full-game signatures $h=0..4$ for Q\_n, even $n=8..24$, in both e.p. heap modes; and P8 as a non-\textasteriskcentered{}3 control.

\textbf{Changed:} replaced the “by inspection” window claim with explicit windows and actual-node counts; removed the unnecessary symmetric-arc assertion; corrected Q8's definition; separated reduced-graph evidence from unreduced signatures and marked B's limited range; stated H1–H8 and the exact meaning of $=\star3$; corrected labels and added exact reproduction commands.


\subsection{Note 3: research/ring\_n\_plus\_2\_proof/PROOF\_odd.md (odd $n$)}
\subsection*{$R_n=\star2$ and $Qo_n=\star3$ for every odd $n\ge 11$ (whole-board game): proof note}
\textbf{Summary}

\begin{enumerate}
\item Theorem: in the whole-board game on $C_n$ (race rule "count", e.p. on, an e.p.-capturable pawn is not passed, the game ends at the first passed pawn), for every odd $n\ge11$ the $(n+2)$-pawn positions $R_n=$ W5 B7 (W2 B4)$^{(n-9)/2}$ W2 W3B6 W3 W2B5 B5 W2 B4 and $Qo_n$ (W4 for W5) are $\star2$ resp. $\star3$ (meaning: H7); both fail at $n=9$ (LLLL).
\item The proof is T50's insertion proof (\texttt{research/ring\_family\_proof/proof.tex}, Lemmas 1-3, with the repaired "pattern file" and the written-out Lemma 3 of \texttt{PROOF.md} 2b), at ONE fixed cut $c=6$ and ONE base $n_0=19$ for both families. The wall (files $2..n-7$) has odd length $n-8$ and W2 at both ends; neither end enters the lemmas (H1, H3).
\item C1: the reduced games at $n_0=19$ have 429 ($R$) and 1287 ($Qo$) nodes, no e.p. flag, no passed pawn, zone $\{n-7..n-3,n-1\}$ (for $Qo$ also $\{0,1\}$), signature NNPN resp. NNNP; files $2..11$ are exact pattern files in every node, and $R_{n+2}=\mathrm{ins}_6(R_n)$ was verified for $19\to21,\dots,57\to59$ (A) and, by independent programs, for $19\to21,\dots,31\to33$. C2: $n=11,13,15,17$ by direct search. E.p. convention: a flag is recorded only if an enemy pawn stands beside the landing square (section 1).
\item Status: PROVED [AI] for odd $n\ge19$ given T50's Lemmas 1-3, 2b and C1; CHECKED-2 for $n=11..17$ (finite checks by 5-6 programs each, section 7). By AI models only: two adversarial Claude verifier lenses (section 7) found no mathematical defect (the fixes are wording, convention and labels), but this note has no non-Claude referee yet (the lemma text it relies on, PROOF.md 2b, was refereed by Codex in T53; PROOF\_Q.md by Codex in T62).
\item Scope: whole-board game, race rule "count" only (not cylcgt's component-sum model, not the rules "mover"/"opponent"), odd $n\ge11$; $n=9$ is false.
\end{enumerate}
\subsubsection*{0. Claims and labels}
PROVED = by hand, given the lemmas (all hand proofs here are [AI]). CHECKED-2 = two or more independently written programs agree (the programs are listed in each row; "CHECKED-2" is the label for any number of programs $\ge2$). COMPUTED-1 = one program. Programs: \textbf{A} = T50's \texttt{rg.c} (via \texttt{pump\_verify.py}, \texttt{graph.py}, \texttt{rg solve}); \textbf{B} = \texttt{reduced\_cyl.py}, a rebuild on cylcgt's \texttt{pawns.moves/passed/\_child}; \textbf{E} = \texttt{verify\_compute/myeng.py}, an engine written from the rules (here driven by \texttt{odd/check\_odd.py}, \texttt{odd/base\_cases.py}); and, added by the verifiers of section 7: \textbf{C} = \texttt{check3/eng.py} (T50 verifier's engine, run by \texttt{verify\_computeOdd/run\_b.py}, \texttt{run\_a.py}), \textbf{D} = \texttt{verify\_logicOdd/engineD.py} (a copy of \texttt{verify\_logic/engine.py}, written from the rules by the logic verifier of PROOF.md), \textbf{O} = \texttt{verify\_computeOdd/oddeng.py} (written from the rules for this note by the compute verifier; it uses no code of A, B, E, C or cylcgt). A, C, D, E, O were written from the same rule description (PROOF.md section 0); B uses the project's \texttt{pawns.py}. D and O were differentially tested against \texttt{pawns.py} (S5); E was in PROOF.md section 0.

\par\noindent{\small\setlength{\tabcolsep}{3pt}\begin{tabular}{@{}>{\raggedright\arraybackslash}p{0.074\linewidth}>{\raggedright\arraybackslash}p{0.545\linewidth}>{\raggedright\arraybackslash}p{0.361\linewidth}@{}}
\toprule
id & claim & label \\
\midrule
L1-L3 & T50 Lemmas 1-3 (reduction, wall pawns, insertion) and the written-out proof of Lemma 3 (PROOF.md 2b) & PROVED in T50 / PROOF.md (2b: Codex SOUND, \texttt{ring\_n\_plus\_2\_check/REPORT.md}); cited, not re-proved \\
H1-H7 & hypotheses of the lemmas for $R_n$, $Qo_n$, $c=6$, $n_0=19$; induction; meaning of $=\star k$ & PROVED [AI] given C1 and L1-L3 \\
C1R & $R_{19}$: 429 nodes, 805 K-, 653 T-moves, no flag (canonical e.p. convention, section 1), no passer, zone $\{12..16,18\}$, signature NNPN, files $2..11$ exact pattern in every node & CHECKED-2 (A, B, E; rebuilt by C, D, O in section 7: six programs, node for node) \\
C1Q & $Qo_{19}$: 1287 nodes, 2996 K-, 1959 T-moves, no flag, no passer, zone $\{0,1,12..16,18\}$, NNNP, files $2..11$ exact pattern & CHECKED-2 (A, B, E; rebuilt by C, D, O: six programs) \\
C2 & $R_n+\star h$: NNPN, $Qo_n+\star h$: NNNP ($h\le3$) for $n=11,13,15,17$ (unreduced, no poison pruning, short-circuit solvers; also the reduced-graph solution) & CHECKED-2 (A, E; also C, D, O: five programs); no-early-exit solvers only for $R_{11}$ (D, O), $Qo_{11}$ (O) and $R_{13}$ (O), section 6.3 \\
C3 & $n=9$: signature LLLL (A), LLLLL (E, D, O, C) for both families & CHECKED-2 (A, C, D, E, O) \\
T19 & $R_n=\star2$, $Qo_n=\star3$ for odd $n\ge19$ & PROVED [AI] given L1-L3, C1R/C1Q \\
T & the same for odd $n\ge11$ & PROVED [AI] given C1, C2 \\
S1 & $\mathrm{ins}_6$ pump $19\to59$ (20 steps): start, nodes, K- and T-edges map onto the next graph; 429 / 1287 nodes at every odd $n=19..59$ & COMPUTED-1 (A) for $n>33$; CHECKED-2 for $19\to21\to\dots\to33$ (A, O; D to $27$ ($R$) resp. $25$ ($Qo$); E for $19\to21\to23$) \\
S2 & unreduced signature NNPN / NNNP for every odd $n=11..59$ & COMPUTED-1 (A, \texttt{unreduced\_*\_9\_59.txt}) for $n>25$; CHECKED-2 up to $n=25$ (A, C, E, O; D to $19$); O alone to $n=41$ ($h\le5$); not used beyond $n=17$ \\
S3 & Lemmas 2 and 3 on the actual nodes (section 4); sensitivity controls; D's random tests with arbitrary content outside the arc & CHECKED-2 (E, D, O); sanity only \\
S4 & cylcgt rebuild of the reduced graph identical to A's for $n=11..21$ & CHECKED-2 (A, B) \\
S5 & the rules of D, O (and E, C) agree with \texttt{pawns.py}: moves, passed pawns, promotion distance, verdict with the count race, e.p. & CHECKED-2 (D: 7463 positions, 222,556 child records; O: 64,886 positions, 755,453 moves; 0 differences) \\
S6 & e.p. convention: raw flags (\texttt{pawns.py}: a flag after every double step) give 434 / 1302 nodes at every $n\ge11$, equal to the canonical graphs once inert flags are dropped & CHECKED-2 (O; integrator's own rebuild on \texttt{pawns.py} at $n=19$) \\
\bottomrule
\end{tabular}}\par\smallskip
\subsubsection*{1. Setting}
Model, classes T/L/P/K, reduced game: as in \texttt{proof.tex} and PROOF.md section 1 (race rule "count"). \textbf{E.p. convention} (as in \texttt{proof.tex}, Model): an e.p. flag is recorded only if an enemy pawn stands beside the landing square of the double step; a flag nobody could use is dropped. "No flag" in C1, H4, H6 and Lemmas 2-3 means this canonical convention. A, C, D, E, O use it, and B applies it through \texttt{difftest.norm\_ep}; \texttt{cylcgt/pawns.py} itself records a flag after every double step, but its \texttt{ep\_capturable} and \texttt{moves} act on a flag only when an enemy pawn stands beside the square, so an inert flag changes nothing but the name of the position. This was checked: with raw flags the reduced graphs have 434 ($R$) and 1302 ($Qo$) nodes for every odd $n=11..25$ (5 resp. 15 nodes carry an inert flag; 231 and 639 at $n=9$), and dropping the inert flags gives exactly the canonical graph, edge for edge, with the same signature (O, \texttt{verify\_computeOdd/logs/b\_reduced.txt}; also rebuilt by the integrator (the session that wrote section 7) directly on \texttt{pawns.moves/passed/\_child} at $n=19$: 434 nodes, 5 inert flags, node and edge sets equal to the canonical ones, 1100 raw vs 1082 canonical edges; likewise 1302, 15, 3886 vs 3827 for $Qo_{19}$). Files $0..n-1$, $n$ odd $\ge11$: file 0 W5 ($R$) or W4 ($Qo$); file 1 B7; file $j$ ($2\le j\le n-7$) W2 if $j$ is even, B4 if $j$ is odd; then W3B6 ($n-6$), W3 ($n-5$), W2B5 ($n-4$), B5 ($n-3$), W2 ($n-2$), B4 ($n-1$). That is $2+(n-8)+8=n+2$ pawns. The \textbf{wall} is files $2..n-7$: type E (black pawn two ranks above the white one), length $n-8$ (odd), W2 on file 2 (beside B7) and on file $n-7$ (beside W3B6). \textbf{Pattern file} (PROOF.md section 2): a wall file whose content is exactly its start pawn. $\mathrm{ins}_c$ as in Lemma 3. $R_n$ ($Qo_n$) also denotes the reduced game, $\mathrm{start}_n$ its root. $Z_n=\{n-7,\dots,n-3,n-1\}$ ($R$), $\{0,1\}\cup\{n-7,\dots,n-3,n-1\}$ ($Qo$). Lemma 1: a side wins $G+\star h$ moving first iff it wins the reduced game moving first. Lemma 2: for a wall pawn $q$ of the player $Z$ to move with pattern neighbour files, (a) $Z$ has no passed pawn after any move of $q$, (b) if $c$ has no flag and the window $j-1..j+3$ (or $j-3..j+1$) consists of pattern files, every move of $q$ is L or P. Lemma 3: if $p$ has no passer, no flag, and files $c-4..c+5$ are pattern files, then the K-children and the existence of T-moves of $\mathrm{ins}_c(p)$ are the $\mathrm{ins}_c$-images of those of $p$.

\subsubsection*{2. The cut and the base}
Zone and pattern files. By C1, in every node of the reduced game of $G_n$ ($n\ge19$) all files outside $Z_n$ keep their start content. So files $2..n-8$ are pattern files of every node. The arc $A=\{c-4..c+5\}$ lies in $[2,n-8]$ iff $6\le c\le n-13$; the least $n$ is $n_0=19$, where $c=6$ is the only cut (\texttt{pump\_verify.py} finds $c=6$ automatically at $19\to21$), and $c=6$ is admissible for every $n\ge19$. $A=\{2..11\}$ never wraps: the outside $12..n-1,0,1$ is one cyclic arc, which contains $Z_n$ and the files 0, 1 (not touched by $\mathrm{ins}_6$). $\mathrm{ins}_6$ fixes files $0..6$ and shifts files $\ge7$ by 2, so $\mathrm{ins}_6(Z_n)=Z_{n+2}$. The range $6\le c\le n-13$ is what Lemma 3's arc hypothesis needs, so it is sufficient, not necessary: the computation shows that $\mathrm{ins}_c:G_n\to G_{n+2}$ is an isomorphism of the reduced graphs for every $c=2..n-9$ (O, $n=19..31$), and for $n=11..17$ at some cuts (only $c=2$ at $n=11$; $c=2..8$ at $n=17$; \texttt{iso\_small.txt}). So the graphs for $n<19$ are also $\mathrm{ins}$-images of each other, but they do not satisfy Lemma 3's hypothesis (no ten-file arc of unchanged wall files) and are covered by C2, not by the induction.

\subsubsection*{3. Hypotheses of Lemmas 2 and 3}
\textbf{H1 (wall, phase, ends).} The wall alternates W2 (even) / B4 (odd). $\mathrm{ins}_6$ inserts copies of file 7 (B4) and file 6 (W2) after file 6, so the new files 7, 8 continue the phase; $\mathrm{ins}_6(\mathrm{start}_n)=\mathrm{start}_{n+2}$ for $n\ge14$ (files 6, 7 are wall files), which covers $n\ge19$. The phase and ends differ from $E_n$, $P_n$, $Q_n$ (here W2 at both ends), but Lemma 2(b) uses only the rows W2$\to$3, W2$\to$4 (e.p.) and B4$\to$3 of T50's table, and Lemma 3 only the views of pawns (PROOF.md H2, 2b); neither uses a file's parity or a wall end. \textbf{H2 (no capture by $q$).} With pattern neighbours the capture squares of $q$ on rank 3 are empty, and B4 cannot take e.p. a W2 that has not moved (PROOF.md H3). \textbf{H3 (windows at the two ends of the arc).} Central pawns: files $c-3..c+4=3..10$ of $p$, and $3..12$ of $\hat p=\mathrm{ins}_6(p)$ (arc $2..13$). Window of $q$ on file $j$: $j-1..j+3$ if $j\le c+2$ (resp. $c+4$ in $\hat p$), else $j-3..j+1$; all lie in the arc. \emph{Left end:} $j=3$, $q=$B4, window $2..6$ (file 2 is the first wall file, W2 beside B7 on file 1; files 0, 1 are in no window): after B4$\to$3 the W2 on file 4 takes on $(3,3)$ and is passed (files $2,3,4$ then hold no black pawn), the black B4 on file 5 stays blocked by W2 on file 6. \emph{Right end:} $j=10$ in $p$ ($j=12$ in $\hat p$), $q=$W2, window $7..11$ ($9..13$): after W2$\to$3, resp. $\to$4, B4 on $j-1$ takes (e.p. after $\to$4); the W2 on $j-2$ stays blocked by B4 on $j-3$. The window ends at file $11=n_0-8$ (resp. 13), so the wall-end file $n-7$ (W2, inside $Z_n$) is never part of a window. In 18 of 429 ($R$), 54 of 1287 ($Qo$) nodes the double step of $q$ on file 10 is L, not P: the W2 on file 12 has left, so B4 on file 11 is passed, the "$Y$ has a passer" branch of Lemma 2(b); the proof needs only that central moves are neither T nor K, so nothing breaks (as in PROOF.md H2, end case). \textbf{H4 (start).} $\mathrm{start}_n$ has no passed pawn and no flag: W5/W4 (0) is blocked by B7 (1), B7 by W5/W4; a wall W2 by the B4 on an adjacent file (file 2: B4 on 3; file $n-7$: B4 on $n-8$), a wall B4 on file $j$ by W2 on $j-1$; W3B6 (file $n-6$): each blocks the other; W3 ($n-5$) by B6; W2B5 ($n-4$): each blocks the other; B5 ($n-3$) by W2 ($n-4$); W2 ($n-2$) by B5; B4 ($n-1$) by W2 ($n-2$). E ran the passed-pawn test on all nodes: none has a passer. \textbf{H5 (race rule).} Lemma 2 never uses the race; in Lemma 3 the passed pawns of $c_\mu$ and $\mathrm{ins}_6(c_\mu)$ coincide (V3), so any rule that depends on the passed pawns and on who moved gives the same class (PROOF.md H4). The value is rule-specific, the lemmas are not. \textbf{H6 (invariant).} $I_n$: no node of the reduced game has a flag, and every node agrees with $\mathrm{start}_n$ outside $Z_n$. $I_{19}$ is C1. If $I_n$ holds ($n\ge19$): every node is non-terminal with no flag, and $A$ consists of pattern files (section 2), so Lemma 3 applies to every node with $c=6$; with H4 and induction on the distance from the root, nodes, K-edges and T-flags of $R_{n+2}$ are the $\mathrm{ins}_6$-images of those of $R_n$. An image has no flag; it agrees with $\mathrm{start}_{n+2}$ outside $Z_{n+2}$ since files 7, 8 are copies of pattern files 7, 6 and $\mathrm{ins}_6(Z_n)=Z_{n+2}$. So $I_{n+2}$. \textbf{H7 (meaning of $=\star k$; heaps).} As in T50 and PROOF.md H8: $G=\star k$ means that for every heap $\star h$, $h<256$ (a passed pawn is worth $\pm256$), $G+\star h$ and $\star k+\star h$ have the same outcome, namely P for $h=k$ and N otherwise. It is the whole-board/heap-sum convention, not $G+X\equiv\star k+X$ for every game $X$. \emph{R, $k=2$:} the signature of $R_n+\star h$ is N, N, P for $h=0,1,2$ (C1R); for every $h\ge3$, $R_n+\star h$ has the heap move to the P-position $R_n+\star2$, so it is N. \emph{Qo, $k=3$:} N, N, N, P for $h=0..3$; for $h\ge4$ the move to $Qo_n+\star3$ gives N. The start has no flag and neither do the reduced nodes, so the e.p. convention across a heap move is immaterial (E ran both readings "keep" and "expire": identical results, \texttt{base\_cases.log}).

\subsubsection*{4. Finite checks (rerun: \texttt{sh research/ring\_n\_plus\_2\_proof/odd/run\_odd.sh}; logs in \texttt{logs/} and \texttt{odd/})}
\textbf{C1R, C1Q ($n_0=19$).} A: \texttt{logs/pump\_R\_fixc.txt}, \texttt{pump\_Qo\_fixc.txt} (\texttt{FIXC=6 python3 pump\_verify.py R 19 59}, same for Qo): $n=19$ has 429 nodes, 805 K-, 653 T-moves, zone $\{12,13,14,15,16,18\}$, ep=0, NNPNNNNN ($h\le7$), resp. 1287, 2996, 1959, zone $\{0,1,12..16,18\}$, NNNPNNNN. B: \texttt{logs/reduced\_cyl\_R.txt}, \texttt{reduced\_cyl\_Qo.txt}: the cylcgt rebuild equals A's graph (nodes and $(\text{src},\text{side},\text{target})$ triples: 1082 = 805 + 277, resp. 3827 = 2996 + 831) for $n=11..21$, $n=19$ included. E: \texttt{odd/check\_odd\_R.log}, \texttt{check\_odd\_Qo.log}: the same counts, zone, no flag, signature NNPNN / NNNPN ($h\le4$), no passed pawn of either colour in any node, and files $2..11$ (and $2..13$ at $n=21$, $2..15$ at $n=23$) equal to the exact pattern W2 (even), B4 (odd) in every node. C, D, O (section 7) rebuilt both graphs from the rules: the same counts, zone, no flag, no passed pawn of either colour in any node, files $2..n-8$ exact pattern in every node, signatures NNPN / NNNP, for every odd $n=11..25$ (D: $n=11..19$ and $R_{21}$; C: $n\le21$). The integrator reran D's and O's builds at $n=19$ (both families) and got the same numbers. \textbf{Pump (S1).} A: for $n=19,21,\dots,59$ both families keep 429 / 1287 nodes, zone $Z_n$, ep=0, the signature; each of the 20 steps $19\to21,\dots,57\to59$ with the fixed cut $c=6$ reports "start-\textgreater{}start True, nodes onto True, kept+T edges equal True". E: $\mathrm{ins}_6(\mathrm{start}_n)=\mathrm{start}_{n+2}$, $R_{n+2}=\mathrm{ins}_6(R_n)$ node for node with equal K- and T-children for $n=19,21$ (858 resp. 2574 (node, side) pairs per step). D (cut $c=6$): $R$ for $19\to21\to23\to25\to27$, $Qo$ for $19\to21\to23\to25$, all "start-\textgreater{}start, nodes onto, K-edges and T-flags equal" (\texttt{verify\_logicOdd/logs\_pump\_c6.txt}); O: both families for $19\to21\to\dots\to33$ (\texttt{verify\_computeOdd/logs/c\_R.txt}, \texttt{c\_Qo.txt}); the integrator reran D's $R$ step $19\to21$ at $c=6$ (True). The graphs of O, E, C are identical node for node and edge for edge for $n=9..25$ (C for $n\le21$) and equal every number in A's logs \texttt{pump\_R\_11\_59.txt}, \texttt{pump\_Qo\_11\_59.txt} (O's \texttt{run\_b.py}); D's graphs agree with them in every count (nodes, K, T, flags, passers), zone and signature for $n=11..19$ (and $R_{21}$). \textbf{Lemmas 2, 3 on the actual nodes (S3, E).} Lemma 2 on every central wall pawn (files $3..10$ at $n=19$, $3..12$ at $n=21$), both colours, with the windows of H3: $R$: 3432 q-positions, 5148 moves at $n=19$ (4290, 6435 at $n=21$); $Qo$: 10296, 15444 ($n=19$), 12870, 19305 ($n=21$); (a) 0 failures, (b) 0 T or K moves, every move L or P (file 3 and the left end: all P; right end $j=c+4$: 18 L + 840 P for $R$, 54 L + 2520 P for $Qo$). Lemma 3 recomputed on every $\mathrm{ins}_6(p)$ (not via the explored graph): 0 mismatches. Control: for the cuts $c=7..10$ Lemma 3's conclusion still holds on the nodes of $n=19$ although the arc meets the zone (radius 5 is conservative), for $c=11$ (arc $7..16$) it fails at 48 of 858 ($R$) and 150 of 2574 ($Qo$) (node, side) pairs, so the test can fail. That control is weak (it only looks at reachable nodes). A stronger one is D's random test with arbitrary content outside the arc (\texttt{verify\_logicOdd/t\_random.py}): positions are the start with 1-4 files outside the arc $2..11$ replaced by random contents, or a node plus 1-2 arbitrary moves, kept if they have no passer and no flag; K-children and the T-flag of $p$ are compared with those of $\mathrm{ins}_6(p)$ (about 100k K-children compared in all, 0 mismatches: $R_{19}$ 4000, $Qo_{19}$ 4000, $R_{21}$ 2000, $Qo_{25}$ 1500, $R_{27}$ 1500 positions, plus 2500 + 2500 node-walk positions). Cutting at $c\ne6$ while the same outside content is randomized gives 0 mismatches at $c=5,7,9$ (the radius-5 hypothesis is conservative) and mismatches at $c=10$ (208 of 1500), $c=11$ (233) and $c=12$ (1311), so the test fails when the hypothesis fails (\texttt{logs\_random\_ctrl.txt}; the integrator's rerun with other seeds: 0 mismatches at $c=6$ and $c=7$, 51 of 400 at $c=10$). Lemma 2 with garbage outside the five-file window (two runs of 30000 positions per family, every wall file, both window directions): 0 violations of (a) or (b). On the nodes, D compares every move of every non-central pawn with its $\mathrm{ins}_6$-image (class T/L/P/K): $R_{19}$ 7131 moves, $Qo_{19}$ 20964, 0 mismatches; 242 resp. 726 of these are races (both sides passed), including ties (1 resp. 3), and all agree. O recomputes Lemma 3 on every $\mathrm{ins}_6(p)$ at $n=19,21,23$ (858 resp. 2574 (node, side) pairs per $n$): 0 mismatches. D and O reproduce the Lemma 2 tallies above exactly, including the right-end double step (18 of 429 and 54 of 1287 nodes: L because $q$'s double step leaves B4 on file 11 passed once W2 on file 12 has moved), and O repeats Lemma 2 with every wall file $3..n-8$ having an exact-pattern window inside the files $2..n-7$ (0 failures). \textbf{C2, C3 (base cases).} Direct search, no poison pruning: A \texttt{logs/unreduced\_R\_9\_59.txt} (odd $n=9..59$: LLLL at 9, then 25 lines NNPN; nodes 2916, 3654, 4392, 5130 for $n=11,13,15,17$) and \texttt{unreduced\_Qo\_9\_59.txt} (LLLL, then 25 lines NNNP; 7780, 9749, 11718, 13687). E \texttt{odd/base\_cases.log}, $h\le4$, "keep" and "expire": NNPNN ($R$) / NNNPN ($Qo$) for $n=11..21$ (4628, 5832, 7036, 8240 states for $R$, $n=11..17$; 12239, 15451, 18663, 21875 for $Qo$), LLLLL at $n=9$. Resources: each job one process under \texttt{nice -{}n 15}; the Python jobs peaked at 27 MB (E) and 45 MB (A's pump, \texttt{time -{}l}). \emph{Further runs (section 7).} The same short-circuit search, $h\le6$, in O, E and C for every odd $n=9..25$: NNPNNNN ($R$, $n\ge11$) / NNNPNNN ($Qo$), LLLLLLL at $n=9$, with identical state counts and both e.p.-across-heap readings (\texttt{verify\_computeOdd/logs/a\_*.txt}); in D for $n=9..19$ (\texttt{verify\_logicOdd/logs\_unreduced\_*.txt}), each time also against the signature of the reduced graph, which is a second, different route to C2 (Lemma 1). Searches that evaluate every child and every heap move (no early exit, $h\le4$): $R_9$ and $Qo_9$ LLLLL (D, O); $R_{11}$ NNPNN (D: 885,125 states, O: 885,125 keep / 836,505 expire; the two engines give the same state count); $Qo_{11}$ NNNPN (O: 1,662,552 / 1,567,783 states); $R_{13}$ NNPNN (O, keep only: 4,261,697 states). Not run without early exit: $R_{15}$, $R_{17}$, $Qo_{13}$, $Qo_{15}$, $Qo_{17}$ (memory), so for these the short-circuit solvers (A, C, D, E, O, which agree) and the reduced-graph solution carry C2. O also gives $R_n=\star2$, $Qo_n=\star3$ ($h\le5$) for odd $n=27..41$, which is not used.

\subsubsection*{5. Proof of the Theorem}
Let $n\ge19$ be odd. $I_{19}$ is C1; H6 gives $I_n$ for all odd $n\ge19$ and $R_n\cong R_{19}$, $Qo_n\cong Qo_{19}$ (nodes, K-edges, T-flags). By C1 and Lemma 1, $R_n+\star2$ and $Qo_n+\star3$ are P-positions, and the signatures are NNP.. resp. NNNP.. as in C1; H7 gives all heaps: $R_n=\star2$, $Qo_n=\star3$. For $n=11,13,15,17$ use C2 and H7 (the same heap step). $\blacksquare$ (T19: PROVED [AI] given L1-L3, C1; T: with C2.)

\subsubsection*{6. Gaps and caveats}
\begin{enumerate}
\item \textbf{No gap found in the hypotheses.} What is new relative to PROOF.md and PROOF\_Q.md is only the wall's phase and ends (W2 at both ends, odd length), the cut $c=6$ at the left end of the pattern run (the window of the pawn on file 3 starts at the first wall file 2, beside B7) and $n_0=19$; H1-H3 treat each. The one point to look at is the window of the pawn on file 3: it needs only that file 2 is a pattern file, which holds in every node (file 2 is outside $Z_n$ and its W2 never moves in the reduced game: its 2 moves per node are P in all nodes, C1 information). The logic verifier of section 7 re-derived by hand every point where the odd configuration differs from the even cases (left end $j=3$ with window $2..6$, right end $j=10$ in $p$ and $j=12$ in $\hat p$ with the e.p. double step, odd wall length, W2 at both ends, the zones touching the arc, e.p. and race handling) and changed no step; the compute verifier recomputed the facts these arguments cite (zones, windows, tallies, pump).
\item \textbf{Lemmas 1-3 are cited, not re-proved}; their hand proofs are T50's and 2b, refereed SOUND for $P_n$ (T53, Codex) and $Q_n$ (T62, Codex). This note itself has been through two adversarial Claude verifier lenses (section 7), which is evidence, not a formal referee report; it has no non-Claude referee (suggested: a model other than Claude, with H1, H3, H6 and C1). The truth of Lemma 2(b) itself rests on T50 and the Codex referee of PROOF.md and PROOF\_Q.md.
\item \textbf{C1, C2 are computer facts}: the theorem for $n\ge19$ is "proved modulo C1", the theorem for $n=11..17$ is a finite computation. A, C, D, E, O were written from one rule description (B uses the project's \texttt{pawns.py}); a shared misreading is excluded only by differential tests against \texttt{pawns.py}: D on $R_{19}$, $Qo_{19}$ (7463 positions, 222,556 child records), O on 64,886 positions with 755,453 moves (including e.p. positions and both-sides-passed races), E in PROOF.md (about 61k positions): 0 differences everywhere. The searches for $n=11..17$ are short-circuit solvers; the no-early-exit solver was run for $R_9$, $Qo_9$, $R_{11}$, $Qo_{11}$ and $R_{13}$ ($h\le4$) only, not for $R_{15}$, $R_{17}$, $Qo_{13}$, $Qo_{15}$, $Qo_{17}$ (section 4); for these the short-circuit solvers (five programs) and the solution of the reduced graph (a different route, via Lemma 1) agree. Wherever a no-early-exit run exists it agrees with the short-circuit result, at $n=9$ as well.
\item \textbf{Scope}: whole-board game, race rule "count", e.p. flags recorded only if an enemy pawn stands beside the landing square (game-equivalent to \texttt{pawns.py}'s raw flags, section 1); race-rule sensitivity and the component-sum model were not examined for the odd families; $R_9$, $Qo_9$ are not $\star2$, $\star3$ (C3).
\item The graphs already have 429 / 1287 nodes at $n=11..17$ (zones as in \texttt{logs/pump\_*\_11\_59.txt}), but $n<19$ has no ten-file arc of unchanged wall files, so those $n$ are not instances of the induction and rest on C2. $n=9$ is not an instance of the pump either: its reduced graphs are different (2082 nodes for $R$ with files $0$ and $8$ in the zone, 5276 for $Qo$) and have signature LLLL, as C3 says.
\item The note's sensitivity control on the nodes (Lemma 3 fails only at $c=11$) is weak; the random-outside control of section 4 is the one to cite.
\item Do not claim: the component-sum model, another race rule, even $n$ (that is PROOF.md, PROOF\_Q.md), or equivalence against arbitrary addends.
\end{enumerate}
\subsubsection*{7. Verification}
Two adversarial verifiers judged this note in separate lenses (both are Claude models; neither is a formal referee). Every problem they raised was then re-checked against this file, PROOF.md, \texttt{proof.tex}, PROOF\_Q.md, their scripts and logs, and the edits above were made. Scripts and logs: \texttt{verify\_logicOdd/} (engine D; logs \texttt{logs\_*.txt}), \texttt{verify\_computeOdd/} (engine O; logs in \texttt{logs/}, \texttt{sh research/ring\_n\_plus\_2\_proof/verify\_computeOdd/run\_all.sh} from the project root, one job at a time, \texttt{nice -{}n 15}; the no-early-exit $R_{13}$ run is \texttt{exh13.py}, not in \texttt{run\_all.sh}). D's scripts are run from \texttt{verify\_logicOdd/} with \texttt{PYTHONPATH=} the project root: \texttt{t\_graph.py FAM n...}, \texttt{t\_pump.py FAM c n\_lo n\_hi}, \texttt{t\_random.py MODE FAM n N seed [c]} (modes \texttt{l3mut}, \texttt{l3walk}, \texttt{l2gar}, \texttt{ctrl}), \texttt{t\_nodes.py FAM n}, \texttt{t\_unreduced\_odd.py FAM n...}, \texttt{t\_diff.py} (differential test against \texttt{pawns.py}; run from the project root).

\textbf{Verdicts.} Logic lens: ACCEPT (could not refute that T50's Lemmas 1-3 transfer to $R_n$, $Qo_n$ at $c=6$, $n_0=19$; five non-blocking caveats). Compute lens: ACCEPT (no difference from the note in (a) the unreduced signatures, (b) the reduced graphs, (c) the pump and the Lemma 2, 3 tallies; four clarifications; it did not examine the hand proofs of Lemmas 1-3).

\textbf{What the logic verifier recomputed} (D, a copy of \texttt{verify\_logic/engine.py}; independent of A, B, E). C1: $R_{19}$ 429 nodes, 805 K, 653 T; $Qo_{19}$ 1287, 2996, 1959; no flagged node, no passer in any node, zones $\{12..16,18\}$ and $\{0,1,12..16,18\}$, files $2..11$ exact pattern in every node, signatures NNPN / NNNP; the same at $R_{21}$ and at $n=11..17$ ($R$) and $n=11,13$ ($Qo$). Pump at $c=6$: $R$ for $19\to21,\dots,25\to27$, $Qo$ for $19\to21,\dots,23\to25$, all true. Differential test against \texttt{cylcgt/pawns.py} (moves, passed pawns, terminal verdict with the count race): $R_{19}$ 4209 positions, 127,244 child records; $Qo_{19}$ 3254 positions, 95,312 child records; 0 mismatches. Random stress tests of Lemma 3 with arbitrary content outside the arc (about 100k K-children, 0 mismatches) and of Lemma 2 with garbage outside the window (two runs of 30000 positions per family, 0 violations); the corrected control fails at $c=10,11,12$ and passes at $c=5,7,9$. Lemma 2 (a), (b) on every central wall pawn of every node at $n=19,21$ (the tallies of section 4, including the 18 / 54 L double steps at the right end, which it traced to B4 on file 11 becoming passed); every far-pawn move compared with its $\mathrm{ins}_6$-image, including 242 / 726 races with ties. Unreduced search $R_n$, $Qo_n$, $n=9..19$, both e.p. readings, against the reduced-graph solution; no-early-exit for $R_9$, $R_{11}$, $Qo_9$. It re-derived by hand the left end ($j=3$, window $2..6$), the right end ($j=10$ in $p$, $j=12$ in $\hat p$) and the arc/zone arithmetic, and found them as stated in H1-H3.

\textbf{What the compute verifier recomputed} (O = \texttt{oddeng.py}, written from the rules, stdlib only; plus E and C with its own drivers; it used none of A's \texttt{rg.c}, cylcgt's game-value code, \texttt{check\_odd.py} or \texttt{base\_cases.py}). (a) $G+\star h$, $h=0..6$, odd $n=9..25$ in O, E, C: LLLLLLL at $n=9$, NNPNNNN ($R$) / NNNPNNN ($Qo$) for every odd $n=11..25$, identical state counts, both e.p.-across-heap readings and raw flags; $h\le5$ for odd $n=27..41$ in O. No-early-exit: $R_9$, $Qo_9$, $R_{11}$, $Qo_{11}$ (both readings, $h\le4$) and $R_{13}$ (keep). (b) Reduced graphs for odd $n=9..25$: 429 / 805 / 653 and 1287 / 2996 / 1959 for every $n\ge11$, 277 and 831 (node, side) pairs with a T-move, node for node and edge for edge identical to E ($n=9..25$) and C ($n=9..21$), every number equal to A's \texttt{pump\_*\_11\_59.txt}; at $n=9$ the graphs are different (2082, 5276 nodes) and the signature is LLLL. (c) $\mathrm{ins}_6$: start to start and nodes onto nodes with K- and T-children preserved for $19\to21,\dots,31\to33$ in both families; Lemma 3 recomputed on every $\mathrm{ins}_6(p)$ at $n=19,21,23$: 0 mismatches; Lemma 2 tallies equal to the note's; control: 0 failures at $c=7..10$, 48 of 858 ($R$) and 150 of 2574 ($Qo$) at $c=11$. The isomorphism holds for every cut $c=2..n-9$. Raw e.p. flags: 434 / 1302 nodes, game-equivalent to the canonical graphs. Differential test of O against E, C and \texttt{pawns.py} (rules only): 64,886 positions (seeds 1, 7), 755,453 moves, 0 differences in moves, passed pawns, promotion distance or verdict.

\textbf{My own re-check} (integrator; spot runs under \texttt{nice -{}n 15}, at most two processes at a time, no git, no kill; scratch script not saved in the project). D rebuilt $R_{19}$ (429 / 805 / 653, zone $\{12..16,18\}$, NNPNN) and $Qo_{19}$ (1287 / 2996 / 1959, zone $\{0,1,12..16,18\}$, NNNPN); O's \texttt{run\_b.py} at $n=19$ for both families gave the same numbers, NNPNNNNN / NNNPNNNN, E and C graph-identical, 434 / 1302 raw-flag nodes equivalent to the canonical ones. I rebuilt the raw-flag graph directly on \texttt{pawns.moves/passed/\_child} (no \texttt{norm\_ep}, so a flag after every double step): $R_{19}$ 434 nodes, 5 with a flag, none live (an enemy pawn beside the landing square), 1100 edges; $Qo_{19}$ 1302, 15, 3886; after dropping inert flags the node and edge sets equal the canonical ones (1082, 3827 edges) in both. From \texttt{pawns.py} itself: \texttt{moves} sets a flag after every double step, \texttt{ep\_capturable} and the e.p. capture act only if an enemy pawn stands beside the square, and \texttt{proof.tex} (Model) says a flag is recorded only in that case, which is the convention now stated in section 1. D's tests with new seeds: Lemma 3 at $c=6$ on 400 random positions, 0 mismatches; the same test cut at $c=7$: 0, at $c=10$: 51 of 400 mismatches; Lemma 2 with garbage outside the window, $Qo_{19}$, 2000 positions: every move L or P; D's pump $R$ $19\to21$ at $c=6$: true. I checked the numbers this note cites against the files they come from (\texttt{pump\_*\_fixc.txt}, \texttt{reduced\_cyl\_*.txt}, \texttt{unreduced\_*\_9\_59.txt}, \texttt{odd/check\_odd\_*.log}, \texttt{odd/base\_cases.log}): all agree. \texttt{verify\_logicOdd/engineD.py} is byte-identical to \texttt{verify\_logic/engine.py}, and \texttt{oddeng.py} imports only the standard library, as claimed. "Codex refereed 2b" is right: \texttt{delegate/run.sh} and \texttt{delegate/REVIEW.md} give T53 as Codex, though \texttt{REPORT.md} itself does not name the model. Not mine to edit: PROOF.md section 6.2 still says 2b has not been refereed by a second model, which T53 has since done.

\textbf{Problems raised, and what was done.}

\par\noindent{\small\setlength{\tabcolsep}{3pt}\begin{tabular}{@{}>{\raggedright\arraybackslash}p{0.072\linewidth}>{\raggedright\arraybackslash}p{0.390\linewidth}>{\raggedright\arraybackslash}p{0.087\linewidth}>{\raggedright\arraybackslash}p{0.418\linewidth}@{}}
\toprule
\# & problem (source) & confirmed? & fix \\
\midrule
1 & this note has no referee other than Claude models (logic) & yes (non-blocking; was already stated) & summary 4 and 6.2: two Claude verifier lenses have run, there is no non-Claude referee of this note, 2b was refereed by Codex (T53) \\
2 & Lemmas 1-3 are cited, not re-proved; 2b's truth rests on the Codex referee (logic, compute iii) & yes (by design) & 6.2 sharpened; 6.1 records that the logic lens re-derived the odd-specific points by hand and found nothing to change (the compute lens checked the cited facts only) \\
3 & C1, C2 are computer facts; only differential tests guard against a rule misreading shared by the programs (logic, compute i) & yes & 6.3 and table S5 give the differential-test numbers for D, O (and E); "proved modulo C1" kept \\
4 & the note's own sensitivity control (Lemma 3 fails only at $c=11$ on the nodes) is weak (logic) & yes & section 4 cites D's random-outside control (fails at $c=10,11,12$, passes at $c=5,7,9$) and my rerun; 6.6 \\
5 & $n=11..17$ rest on C2; no-early-exit run only for $R_9$, $R_{11}$, $Qo_9$ (logic) and not for $n=15,17$, $Qo_{13}$ (compute ii) & yes (caveat) & exact coverage in the C2 row, section 4 and 6.3: no-early-exit for $R_9$, $R_{11}$, $R_{13}$, $Qo_9$, $Qo_{11}$; the reduced-graph solution and five short-circuit programs for the rest \\
6 & "no flag" holds only under the canonical e.p. convention; \texttt{pawns.py}'s raw flags give 434 / 1302 nodes (compute) & yes (re-derived from \texttt{pawns.py}, \texttt{proof.tex}, and by my own rebuild) & convention stated in section 1 next to C1, in summary 3, 6.4 and table S6: the raw games are equivalent \\
7 & the range $6\le c\le n-13$ is sufficient, not necessary (compute) & yes (information) & one sentence in section 2; $n<19$ stays outside the induction \\
8 & $n=9$ is not an instance of the pump (different graphs) (compute) & yes (consistent with C3) & 6.5 \\
9 & labels: C1 has six programs, C2 five, S1-S3 more than before; "CHECKED-2" should say it means $\ge2$ (both) & yes & section 0 and the table updated with programs, ranges and the new rows S5, S6 \\
\bottomrule
\end{tabular}}\par\smallskip
No problem was rejected. None changes the mathematics: problems 1, 2, 3, 5, 8 are caveats that were already stated and are now exact, problem 4 swaps in a stronger control, problem 6 states a convention the proofs already used, problem 7 is information, problem 9 is labels. There was no new mathematics to referee: the note's lemma transfer is H1-H3, H6, which the logic lens re-derived by hand (the compute lens checked only the finite facts those arguments cite).

\textbf{Final verdict: ACCEPT\_WITH\_FIXES (fixes applied).} Reasons: (1) the logic lens could not refute the transfer of T50's Lemmas 1-3 to the odd wall (hand re-derivation of the left end, right end and arc/zone arithmetic; Lemma 3 stress-tested with arbitrary outside content, with a control that fails at $c=10,11,12$), and the compute lens found no difference in any computed fact the hand argument cites; (2) the finite facts C1, C2, C3 and the pump were rebuilt by three further programs (C, D, O) and my spot runs, node for node, and D and O match \texttt{pawns.py} on 72k positions (about 980k child records) with 0 differences; (3) the remaining weaknesses are named: C1, C2 are finite computations (the theorem for $n\ge19$ is "proved modulo C1"), the lemma text is cited from T50 and PROOF.md 2b, the no-early-exit solver covers only $n\le13$ (and $Qo_{13}$ is missing), no non-Claude referee has read this note, and the result is for the race rule "count" and the whole-board game only.

\textbf{Sentence the paper may use.} In the whole-board game on the cylinder (race rule "count", en passant on, an e.p.-capturable pawn is not passed), the odd wall rings $R_n$ (files W5, B7, a W2/B4 wall of odd length $n-8$ ending in W2, then W3B6, W3, W2B5, B5, W2, B4) and $Qo_n$ (W4 for W5), each of $n+2$ pawns, are $\star2$ resp. $\star3$ for every odd $n\ge11$ (and not for $n=9$), meaning that $R_n+\star h$ and $\star2+\star h$ (resp. $Qo_n+\star h$ and $\star3+\star h$) have the same outcome for every nim-heap $\star h$ of size below 256: for $n\ge19$ this follows from T50's insertion argument adapted by hand [AI] and one finite check each (the reduced games at $n=19$, 429 resp. 1287 nodes) [C2], and for $n=11,13,15,17$ and for the failure at $n=9$ from direct finite searches [C2], all confirmed by independently written programs.

Do not add: that this holds in the component-sum model of cylcgt, for another race rule, for even $n$ (PROOF.md and PROOF\_Q.md), or against arbitrary addends rather than nim-heaps; and do not call the hand part refereed by a non-Claude model for the odd case (only the shared lemma text 2b was, by Codex, for $P_n$).


\subsection{Note 4: research/star3\_proof/PROOF\_S3T3.md}
\subsection*{$T_n\equiv\star3$ and $S_n\equiv\star3$ against nim heaps, for every even $n\ge12$ resp. odd $n\ge13$ (whole-board game on $C_n$): proof note}
\textbf{Summary}

\begin{enumerate}
\item Theorem (nim-heap outcome equivalence): in the whole-board game on $C_n$ (race rule "count", e.p. on, an e.p.-capturable pawn is not passed, game ends at the first passed pawn), the $(n+1)$-pawn rings $T_n = (B4\,W2)^{(n-10)/2}\,B5\,W2\,W6B7\,W3\,B7\,B6\,W3\,W2\,B6\,W2$ (even $n\ge12$) and $S_n = B4\,B6\,W3\,(W2\,B4)^{(n-11)/2}\,W6B7\,W3\,B7\,B6\,W3\,W2\,B6\,W2$ (odd $n\ge13$) satisfy: for every heap $h<256$, $G+\star h$ has the outcome of $\star3+\star h$ (P for $h=3$, N otherwise). This is written $G\equiv\star3$. It is not a claim of equality of partizan games in standard CGT (section 3, H7).
\item The proof is T50's insertion proof (\texttt{research/ring\_family\_proof/proof.tex}, Lemmas 1-3, with the written-out Lemma 3 of \texttt{research/ring\_n\_plus\_2\_proof/PROOF.md} 2b), at one fixed cut per family: $T$: $c=4$, $n_0=20$; $S$: $c=7$, $n_0=23$. The $T$ wall is B4/W2 on files $0..n-11$; the $S$ wall is W2/B4 on files $3..n-9$, after the leading $B4\,B6\,W3$. Both are W2/B4 walls in the sense of T50's table; no colour reversal is used.
\item C1: the reduced game at the base has 1305 nodes ($T_{20}$) resp. 40445 nodes ($S_{23}$), no e.p. flag, no passed pawn, zone $\{n-10,n-9\}\cup\{n-7..n-2\}$ resp. $\{0,1,2\}\cup\{n-10,n-9\}\cup\{n-7..n-1\}$, signature NNNPNN ($h=0..5$); files $c-4..c+5$ are exact pattern files in every node; $\mathrm{ins}_c$ maps the graph isomorphically for $n\ge n_0$. C2: every size below $n_0$, i.e. $n=12,14,16,18$ ($T$) and $n=13,15,17,19,21$ ($S$), by direct search. E.p. convention: canonical (a flag only if an enemy pawn stands beside the landing square).
\item Status: PROVED [S] for $n\ge n_0$ given T50's Lemmas 1-3, 2b and the finite checks C1, C2; CHECKED-2 for the base cases. Refereed by two non-Claude models: T71 (Kilo; logic: SOUND WITH FIXES, \texttt{research/star3\_check\_logic/REPORT.md}) and T70 (Codex; SOUND WITH FIXES for the nim-heap claim, with its own independent checker, \texttt{research/star3\_check/REPORT.md}); fixes applied (section 7). Human check pending.
\item Scope: whole-board game, race rule "count" only, nim-heap addends only; $T_{12}$ and $S_{13}$ are the smallest (the $n=10$ ($T$) and $n=11$ ($S$) members are not $\equiv\star3$).
\end{enumerate}
\subsubsection*{0. Claims and labels}
PROVED = by hand, given the lemmas (all hand proofs here are [AI]). CHECKED-2 = two or more independently written programs agree. COMPUTED-1 = one program. Programs: \textbf{A} = T50's \texttt{rg.c} (via \texttt{pump\_verify.py}); \textbf{E} = \texttt{verify\_compute/myeng.py} (written from the rules); \textbf{R} = T70's \texttt{research/star3\_check/independent.py} (written from the rules by Codex, standard library only, imports none of \texttt{cylcgt}, \texttt{rg}, \texttt{myeng.py}). A drives the pump logs; E drives \texttt{check\_S3T3.py} and \texttt{base\_S3T3.py}; R rebuilds the reduced graphs, both insertion steps and every base case.

\par\noindent{\small\setlength{\tabcolsep}{3pt}\begin{tabular}{@{}>{\raggedright\arraybackslash}p{0.074\linewidth}>{\raggedright\arraybackslash}p{0.491\linewidth}>{\raggedright\arraybackslash}p{0.415\linewidth}@{}}
\toprule
id & claim & label \\
\midrule
L1-L3 & T50 Lemmas 1-3 (reduction, wall pawns, insertion) and the written-out Lemma 3 (PROOF.md 2b) & PROVED in T50 / PROOF.md (2b: Codex SOUND, \texttt{ring\_n\_plus\_2\_check/REPORT.md}); cited, not re-proved \\
H1-H7 & hypotheses of the lemmas for $T_n$, $S_n$ & PROVED [AI] given C1 and L1-L3 \\
C1T & $T_{20}$: 1305 nodes, 4018 K-, 1212 T-moves, no flag, no passer, zone $\{10,11,13..18\}$, NNNPNN, files $0..9$ exact pattern in every node; insertion $20\to22$, $22\to24$ exact & CHECKED-2 (A \texttt{logs/pump\_T3\_12\_40.txt}; E \texttt{logs/check\_T3\_20.log}; R \texttt{star3\_check/graphs\_T\_compact.log}) \\
C1S & $S_{23}$: 40445 nodes, 141252 K-, 86822 T-moves, no flag, no passer, zone $\{0,1,2,13,14,16..22\}$, NNNPNN, files $3..12$ exact pattern in every node; insertion $23\to25$, $25\to27$ exact & CHECKED-2 (A \texttt{logs/pump\_S3\_13\_41.txt}; E \texttt{logs/check\_S3\_23.log}; R \texttt{star3\_check/S\_base\_compact.log}, \texttt{S\_step1\_compact.log}, \texttt{S\_step2\_compact.log}) \\
C2 & $T_n+\star h$: NNNPNN for $n=12,14,16,18$; $S_n+\star h$: NNNPNN for $n=13,15,17,19,21$ (both e.p. readings) & CHECKED-2 for every $n$: A (reduced game, pump logs) and R (unreduced direct search, \texttt{star3\_check/bases\_T\_compact.log}, \texttt{bases\_S\_compact\_all.log}); E (unreduced direct search) too for $n=12,14,16$ and $13,15$ \\
C3 & pump: $T$ nodes/zone/signature stable for $n=12..40$; $S$ stable for $n=19..41$ (signature NNNPNN for every $n=13..41$) & COMPUTED-1 (A); the steps used by the proof are in C1 \\
T & $T_n\equiv\star3$ (even $n\ge12$), $S_n\equiv\star3$ (odd $n\ge13$) & PROVED [S] given L1-L3, C1, C2 \\
\bottomrule
\end{tabular}}\par\smallskip
\subsubsection*{1. Setting}
Model, classes T/L/P/K, reduced game: as in \texttt{proof.tex} and PROOF.md section 1 (race rule "count"). A \textbf{pattern file} is a wall file holding exactly its starting pawn and nothing else. E.p. convention: a flag is recorded only if an enemy pawn stands beside the landing square of the double step; a flag nobody could use is dropped. "No flag" in C1, H4, H6 means this canonical convention.

Files $0..n-1$, $n$ even $\ge12$ ($T$) or odd $\ge13$ ($S$):

\begin{itemize}
\item $T_n$: $(B4\,W2)^{(n-10)/2}$ on files $0..n-11$, then $B5$ ($n-10$), $W2$ ($n-9$), $W6B7$ ($n-8$), $W3$ ($n-7$), $B7$ ($n-6$), $B6$ ($n-5$), $W3$ ($n-4$), $W2$ ($n-3$), $B6$ ($n-2$), $W2$ ($n-1$). The tail occupies 10 files and holds 11 pawns ($W6B7$ is two pawns), so the ring has $(n-10)+11=n+1$ pawns. \textbf{Wall} = files $0..n-11$ (length $n-10$, even), $B4$ on even files, $W2$ on odd files, $B4$ at the left end and $W2$ at the right end (file $n-11$ is odd since $n$ is even).
\item $S_n$: $B4$ ($0$), $B6$ ($1$), $W3$ ($2$), then $(W2\,B4)^{(n-11)/2}$ on files $3..n-9$, then $W6B7$ ($n-8$), $W3$ ($n-7$), $B7$ ($n-6$), $B6$ ($n-5$), $W3$ ($n-4$), $W2$ ($n-3$), $B6$ ($n-2$), $W2$ ($n-1$). The tail occupies 8 files and holds 9 pawns ($W6B7$ is two pawns), so the ring has $3+(n-11)+9=n+1$ pawns. \textbf{Wall} = files $3..n-9$ (length $n-11$, even), $W2$ on odd files, $B4$ on even files.
\end{itemize}
Zone $Z_n$ (the files that change in some node of the reduced game): $Z_n=\{n-10,n-9\}\cup\{n-7..n-2\}$ ($T$), $Z_n=\{0,1,2\}\cup\{n-10,n-9\}\cup\{n-7..n-1\}$ ($S$). For $T$ the zone lies outside the wall. For $S$ the zone contains the leading files $0,1,2$ and the last two wall files $n-10,n-9$; the \textbf{fixed run} of $S$ is the wall minus the zone, files $3..n-11$.

\subsubsection*{2. The cut and the base}
By C1, in every node of the reduced game of $G_n$ ($n\ge n_0$) all files outside $Z_n$ keep their start content.

\begin{itemize}
\item $T$: the whole wall $0..n-11$ is pattern in every node. With $c=4$ the arc $A=\{c-4..c+5\}=\{0..9\}$ lies in the wall iff $n-11\ge9$, i.e. $n\ge20$; it does not fit at $n=18$. Set $n_0=20$. $\mathrm{ins}_4$ fixes files $0..4$ and shifts files $\ge5$ by 2, so $\mathrm{ins}_4(Z_n)=Z_{n+2}$, and $c=4$ is admissible for every even $n\ge20$. Pump: insertion isomorphism $20\to22$ at $c=4$.
\item $S$: not every wall file is pattern (the wall-end files $n-10,n-9$ are in the zone); the invariant used is that the fixed arc is pattern. With $c=7$ the arc $A=\{3..12\}$ lies in the fixed run $3..n-11$ iff $n-11\ge12$, i.e. $n\ge23$; at $n=21$ the zone files $11,12$ overlap the arc. Set $n_0=23$. $\mathrm{ins}_7$ fixes files $0..7$ and shifts files $\ge8$ by 2, so $\mathrm{ins}_7(Z_n)=Z_{n+2}$ (at $S_{25}$: $\{0,1,2,15,16,18..24\}$), and the fixed run grows to $3..n-9$ of $S_{n+2}$; the arc stays disjoint from the zone for every odd $n\ge23$. Insertion isomorphisms $23\to25$ and $25\to27$ at $c=7$ are checked by E and R; the pump line for $n=25$ in \texttt{pump\_S3\_13\_41.txt} is its check of $23\to25$ against the previous size, at $c=7$.
\end{itemize}
The ranges are sufficient, not necessary: E's control cuts give exact insertion for $c=5..8$ ($T$, $n=20$) and $c=8..11$ ($S$, $n=23$), first failure at $c=9$ ($T$) and $c=12$ ($S$).

\subsubsection*{3. Hypotheses of Lemmas 2 and 3}
\textbf{H1 (wall, phase, ends).} $T$: the wall alternates $B4$ (even) / $W2$ (odd), a W2/B4 wall as in T50's table; $\mathrm{ins}_4$ inserts after file 4 copies of files 5 ($W2$) and 4 ($B4$), so the new files 5, 6 continue the phase; $\mathrm{ins}_4(\mathrm{start}_n)=\mathrm{start}_{n+2}$ for $n\ge20$. $S$: the wall alternates $W2$ (odd) / $B4$ (even); $\mathrm{ins}_7$ inserts $B4$, $W2$ on the new files 8, 9, continuing the phase; $\mathrm{ins}_7(\mathrm{start}_n)=\mathrm{start}_{n+2}$ for $n\ge23$. The leading files $0,1,2$ of $S$ are in the zone, outside the arc, and fixed by $\mathrm{ins}_7$. Both walls match T50's table directly (Lemma 2 only needs pattern files $j-1,j,j+1$ of alternating colour two ranks apart, in either phase), so no colour reversal is needed.

\textbf{H2 (no capture by $q$).} On a pattern triple, the capture squares of a $W2$ are on rank 3 while the adjacent $B4$ stand on rank 4, and the capture squares of a $B4$ are on rank 3 while the adjacent $W2$ stand on rank 2; so a wall pawn only pushes, and a $B4$ cannot take e.p. a $W2$ that has not moved (PROOF.md H3). After a push an enemy pawn on an adjacent file is still ahead of it; a $W2$ double step to rank 4 beside a $B4$ leaves a live e.p. right and is not passed.

\textbf{H3 (windows at the two ends of the arc).} The central pawns are on files $c-3..c+4$; for $j\le c+2$ the window $j-1..j+3$, and for $j\ge c-1$ the window $j-3..j+1$, lies in the arc, which consists of pattern files in every node ($I_n$, H6). Lemma 2 is local to the window, so Lemma 2(b) makes every central move L or P for every $n\ge n_0$.

\begin{itemize}
\item $T$, $c=4$: central pawns files $1..8$ at $n_0=20$, files $1..10$ at $n=22$. Left end $j=1$ ($W2$), window $0..4$; right end $j=8$ ($B4$), window $5..9$ (at $n=22$: $j=10$, window $7..11$). E found 0 T/K moves at both ends in all nodes, every central move P (\texttt{logs/check\_T3\_20.log}).
\item $S$, $c=7$: central pawns files $4..11$ at $n_0=23$, files $4..13$ at $n=25$. Left end $j=4$ ($B4$), window $3..7$; right end $j=11$ ($W2$), window $8..12$ (at $n=25$: $j=13$, window $10..14$). At the right end the pawn has L moves (28219) as well as P moves (52671); Lemma 2(b) allows L, and there are 0 T/K moves (\texttt{logs/check\_S3\_23.log}). The wall-end files $13,14$ of $S_{23}$ (T and K moves occur there) lie outside the arc and outside both end windows.
\end{itemize}
\textbf{H4 (start).} $\mathrm{start}_n$ has no flag, since no pawn has just double-stepped, and no passed pawn. Every wall pawn is blocked by the opposite-colour pawn on a neighbouring wall file. $T$ core: $B5/W2$, the co-located $W6B7$, then the adjacent pairs $W3/B7$, $B6/W3$, $W2/B6$, and the final $W2$. $S$: the leading $B4$, $B6$, $W3$ are blocked by the seam $W2$ on file $n-1$, the $W3$ on file 2 and the wall $B4$ on file 3; the trailing core after $W6B7$ has the same adjacent-opponent pattern as in $T$. E and R ran the passed-pawn test on all reduced nodes: none has a passer.

\textbf{H5 (race rule).} Lemma 2 never uses the race. The insertion shifts files but keeps every pawn's rank and the three-file pattern around the cut, so in Lemma 3 the passed pawns of $c_\mu$ and $\mathrm{ins}_c(c_\mu)$ and their promotion distances coincide; any rule that depends on the passed pawns and on who moved, "count" included, gives the same class (PROOF.md H4). The value is rule-specific, the lemmas are not.

\textbf{H6 (invariant).} $I_n$: no node of the reduced game has a flag, and every node agrees with $\mathrm{start}_n$ outside $Z_n$; in particular the arc $A$ is pattern in every node ($T$: the whole wall; $S$: the fixed run $3..n-11$, which contains $A$ for $n\ge23$). $I_{n_0}$ is C1. If $I_n$ holds ($n\ge n_0$): every node is non-terminal with no flag, and $A$ consists of pattern files, so Lemma 3 applies to every node with the chosen $c$; with H4 and induction on the distance from the root, nodes, K-edges and T-flags of $G_{n+2}$ are the $\mathrm{ins}_c$-images of those of $G_n$. An image has no flag; it agrees with $\mathrm{start}_{n+2}$ outside $Z_{n+2}$ since the inserted files are copies of pattern files and $\mathrm{ins}_c(Z_n)=Z_{n+2}$. So $I_{n+2}$. The computed steps ($20\to22\to24$; $23\to25\to27$) agree with this correspondence.

\textbf{H7 (meaning of $\equiv\star3$).} As in T50 and PROOF.md H8, and as the theorem is stated: $G\equiv\star3$ means that for every heap $\star h$, $h<256$ (a passed pawn is worth $\pm256$), $G+\star h$ and $\star3+\star h$ have the same outcome: P for $h=3$ and N otherwise. C1 gives signature NNNPNN for $h=0..5$; for $6\le h<256$, the heap move to the P-position $G+\star3$ gives N. The start has no flag and neither do the reduced nodes, so the e.p. convention across a heap move is immaterial (E and R ran both readings "keep" and "expire": identical results). This relation is weaker than equality of partizan games in standard CGT: the numbers 1 and 2 both have outcome L against every nim heap, yet $1\ne2$ (T70). The proof shows neither that the ring games are impartial nor that the nim-heap test determines their standard value, so no standard-value claim is made.

\subsubsection*{4. Finite checks}
\begin{itemize}
\item \textbf{C1T} (\texttt{n\_0=20}): \texttt{PYTHONPATH=. nice -{}n 15 .venv/bin/python research/star3\_proof/check\_S3T3.py T3 20 4} -\textgreater{} \texttt{logs/check\_T3\_20.log}: 1305 nodes, 4018 K-, 1212 T-moves, zone $\{10,11,13..18\}$, NNNPNN, files $0..9$ exact pattern in every node; Lemma 2: 10440 q-positions, 15660 moves, 0 (a)/(b) failures; Lemma 3: 1305 nodes onto 1305 at $20\to22$ and $22\to24$, 0 mismatches; control fails first at $c=9$ (1328 of 2610). A: \texttt{logs/pump\_T3\_12\_40.txt}: 1305 nodes, zone $\{n-10,n-9\}\cup\{n-7..n-2\}$, ep=0, NNNPNNNN for $n=12..40$, insertion $20\to22$ at $c=4$. R: \texttt{star3\_check/graphs\_T\_compact.log}: the same counts, zones and signature at $T_{20},T_{22},T_{24}$, exact insertion $20\to22$ and $22\to24$.
\item \textbf{C1S} (\texttt{n\_0=23}): \texttt{... check\_S3T3.py S3 23 7} -\textgreater{} \texttt{logs/check\_S3\_23.log}: 40445 nodes, 141252 K-, 86822 T-moves, zone $\{0,1,2,13,14,16..22\}$, NNNPNN, files $3..12$ exact pattern; Lemma 2: 323560 q-positions, 485340 moves, 0 failures; Lemma 3: 40445 nodes onto 40445 at $23\to25$ and $25\to27$, 0 mismatches; control fails first at $c=12$. A: \texttt{logs/pump\_S3\_13\_41.txt}: 40445 nodes for $n\ge19$, zone $\{0,1,2\}\cup\{n-10,n-9\}\cup\{n-7..n-1\}$, ep=0, NNNPNNNN for $n=13..41$, insertion $23\to25$ at $c=7$. R: \texttt{star3\_check/S\_base\_compact.log}, \texttt{S\_step1\_compact.log}, \texttt{S\_step2\_compact.log}: the same counts, zone and signature at $S_{23}$; for $23\to25$ and $25\to27$ R recomputed the K/T children at every mapped node for both sides (80890 node/side pairs per step), exact correspondence and closure.
\item \textbf{C2} (every size below $n_0$): E: \texttt{PYTHONPATH=. nice -{}n 15 .venv/bin/python research/star3\_proof/base\_S3T3.py T3 12 16} -\textgreater{} $T_{12},T_{14},T_{16}$ NNNPNN under "keep" and "expire"; \texttt{... base\_S3T3.py S3 13 15} -\textgreater{} $S_{13},S_{15}$ NNNPNN under both readings. R: \texttt{star3\_check/bases\_T\_compact.log} -\textgreater{} $T_{12},T_{14},T_{16},T_{18}$, and \texttt{star3\_check/bases\_S\_compact\_all.log} -\textgreater{} $S_{13},S_{15},S_{17},S_{19},S_{21}$, all NNNPNN under both readings (unreduced whole-board search; 763845 states at $T_{18}$, 15775446 at $S_{21}$). A gives NNNPNN for every one of these $n$ through the reduced game and Lemma 1. So each size below $n_0$ is CHECKED-2 (A and R), and $n\le16$ ($T$), $n\le15$ ($S$) also by E. Reproduce R: the commands in \texttt{star3\_check/REPORT.md}, section "Implementation and reproduction".
\end{itemize}
\subsubsection*{5. Proof of the Theorem}
Let $n\ge n_0$ be of the right parity. $I_{n_0}$ is C1; H6 gives $I_n$ for all $n\ge n_0$ and $G_n\cong G_{n_0}$ (nodes, K-edges, T-flags). By C1 and Lemma 1, $G_n+\star3$ is a P-position and the signature is NNNPNN as in C1; H7 gives all heaps $h<256$: $G_n\equiv\star3$. For the sizes below $n_0$, $T_{12},T_{14},T_{16},T_{18}$ and $S_{13},S_{15},S_{17},S_{19},S_{21}$, use C2 and H7. $\blacksquare$

\subsubsection*{6. Gaps and caveats}
\begin{enumerate}
\item \textbf{No gap found in the hypotheses.} What is new relative to the $M_n/O_n$ proof is the $(n+1)$-pawn tail with the co-located $W6B7$, the leading $B4\,B6\,W3$ of $S$ (in the zone, fixed by $\mathrm{ins}_7$), and the $S$ wall-end files $n-10,n-9$ in the zone, which is why $S$ uses the fixed-arc invariant and $n_0=23$ rather than the arc-in-wall bound $n\ge21$.
\item \textbf{Lemmas 1-3 are cited, not re-proved}; their hand proofs are T50's and 2b, refereed SOUND for $P_n$ (T53, Codex) and $Q_n$ (T62, Codex); 2b was refereed SOUND by Codex (\texttt{ring\_n\_plus\_2\_check/REPORT.md}). T71 flagged 2b as a dependency; it is not re-proved here.
\item \textbf{C1, C2 are computer facts}: the theorem for $n\ge n_0$ is "proved modulo C1"; the sizes below $n_0$ are finite computations. A, E and R agree on every checked position.
\item \textbf{Scope}: whole-board game, race rule "count", canonical e.p. convention; race-rule sensitivity and the component-sum model were not examined for $T_n$, $S_n$; $n<n_0$ are covered by C2, not by the induction.
\item \textbf{The claim is nim-heap outcome equivalence} (H7), not equality to $\star3$ as a standard partizan game. Do not claim: standard CGT equality, the component-sum model, another race rule, or equivalence against arbitrary addends rather than nim-heaps.
\item \textbf{Third program}: the codex \texttt{engine.py} run on the smallest base cases planned in the first draft was not made; R (T70's independent program) replaces it and covers every base case.
\end{enumerate}
\subsubsection*{7. Verification}
Revision (2026-10-29): the fixes of the two referees were applied by the orchestrator. From T70 (Codex, \texttt{research/star3\_check/REPORT.md}): the theorem is stated as nim-heap outcome equivalence, $G+\star h$ has the outcome of $\star3+\star h$ for every $h<256$, not standard CGT equality (Summary 1, H7, section 6.5); the $S$ pawn count now counts 9 pawns in the 8-file tail ($3+(n-11)+9=n+1$), and the $T$ count likewise; the $S$ sentence "the wall files are pattern files" is replaced by the fixed-arc invariant (sections 1, 2, H6); all sizes below $n_0$, including $T_{18}$ and $S_{17},S_{19},S_{21}$, are cited from R's direct search (C2); $n_0=23$ for $S$ is justified and both steps $23\to25$, $25\to27$ are cited, with the pump's $n=25$ line read as its $23\to25$ check. From T71 (Kilo, \texttt{research/star3\_check\_logic/REPORT.md}): all exploratory text is removed; $n_0$ and the cut are stated cleanly for both families (section 2); the $S$ right end in H3 now states that L moves occur, are allowed by Lemma 2(b), and that the classification is $n$-invariant by locality; the coverage gap below $n_0$ is closed (C2) and the C2 label no longer overstates independence. H2 and H4 now carry T70's explicit geometry. The text as T69 wrote it is in \texttt{logs/PROOF\_S3T3\_as\_written\_by\_T69.md}.

Programs and logs as in section 4. The check\_S3T3 runs are E's independent rebuild of the reduced graph and the Lemma 2/3 hypotheses; the base\_S3T3 runs are E's direct search; R rebuilds the reduced graphs, both insertion steps per family and every base case. A, E and R are independently written (C, Python, Python by another model). All pass (logs above). Referees: T71 (Kilo, logic: SOUND WITH FIXES) and T70 (Codex: SOUND WITH FIXES for the nim-heap claim; independent checker \texttt{research/star3\_check/independent.py}).

